\documentclass[11pt,a4paper]{article}
% Preámbulo común a las versiones en inglés y español (compilar con xelatex)
\usepackage{amsmath,amssymb,amsthm,mathtools}
\usepackage{mathpazo}
\usepackage[no-math]{fontspec}
\setmainfont{texgyrepagella}[Path=/usr/share/texmf/fonts/opentype/public/tex-gyre/,Extension=.otf,
  UprightFont=*-regular,BoldFont=*-bold,ItalicFont=*-italic,BoldItalicFont=*-bolditalic]
\usepackage[margin=2.7cm]{geometry}
\usepackage{booktabs,array,graphicx}
\usepackage{enumitem}
\usepackage[round]{natbib}
\usepackage[colorlinks=true,linkcolor=blue!50!black,citecolor=blue!50!black,urlcolor=blue!50!black]{hyperref}
\usepackage{xcolor}
\setlength{\parskip}{0.35em}
\setlist{itemsep=0.2em,topsep=0.3em}
\setlength{\emergencystretch}{3em}

% Notación
\newcommand{\E}{\mathbb{E}}
\newcommand{\PP}{\mathbb{P}}
\newcommand{\ind}{\mathbf{1}}
\newcommand{\X}{\mathcal{X}}
\newcommand{\A}{\mathcal{A}}
\newcommand{\dd}{\,\mathrm{d}}
\newcommand{\Pac}{P_{\mathrm{ac}}}
\newcommand{\Ps}{P_{\mathrm{s}}}
\newcommand{\thac}{\theta_{\mathrm{ac}}}
\newcommand{\ths}{\theta_{\mathrm{s}}}
\newcommand{\Cov}{\operatorname{Cov}}
\newcommand{\Var}{\operatorname{Var}}
\newcommand{\odds}{\operatorname{odds}}
\newcommand{\neff}{n_{\mathrm{eff}}}
\newcommand{\vacio}{\varnothing}

\usepackage[provide=*]{babel}
\babelprovide[main,import]{spanish}
\usepackage{icomma}

\theoremstyle{plain}
\newtheorem{teorema}{Teorema}
\newtheorem{proposicion}{Proposición}
\newtheorem{lema}{Lema}
\newtheorem{corolario}{Corolario}
\theoremstyle{definition}
\newtheorem{ejemplo}{Ejemplo}
\newtheorem{observacion}{Observación}
\renewcommand{\proofname}{Prueba}

\title{Sesgo invisible: cuando los canales miden las covariables de forma distinta}
\author{Pablo Antonio García Pastor\\ \small Investigador independiente \quad \texttt{CORREO}}
\date{Versión en español · 24 de octubre de 2026\\[0.3em]\small La versión oficial es la inglesa; esta es una traducción con el mismo contenido.}

\begin{document}
\maketitle

\begin{abstract}
Las organizaciones quieren saber a menudo qué obtendría un grupo de clientes si lo atendiera otro canal, con solo datos observacionales. La respuesta estándar repondera a las unidades del otro canal, con pesos de un clasificador calibrado. Mostramos cómo puede fallar esa respuesta cuando, como es típico de los datos de canal, cada canal registra las covariables a su manera. Una variable que solo registra un canal rompe el soporte por construcción; el clasificador ideal lo delata, y descartarla es seguro cuando no predice el resultado en el canal básico dadas las demás. Una variable medida con un instrumento distinto en cada canal es peor: el soporte puede sobrevivir, y el estimador converge a un valor desplazado que ningún análisis de los datos observados distingue de la verdad. Derivamos ese desplazamiento exactamente. Con datos públicos en los que existen los dos instrumentos para las mismas personas (índice de masa corporal declarado y medido en NHANES), el desplazamiento derivado coincide de forma aproximada con el sesgo observado, y los diagnósticos habituales de solapamiento mejoran en lugar de delatarlo. Por el camino, leemos el fallo de soporte en el clasificador mediante la descomposición de Lebesgue, mostramos que el estimador doblemente robusto extrapola en silencio donde falta soporte, damos pesos para grupos objetivo definidos por un corte en una puntuación conocida y adaptamos cotas de sensibilidad para la confusión no observada. Una referencia experimental (NSW/CPS) y una simulación cuyas predicciones se registraron antes de ejecutarla (informamos de las que fallaron) completan la evidencia.

\end{abstract}

\section{Introducción}\label{sec:intro}

Las organizaciones reparten a sus usuarios entre canales de calidad y coste distintos: un centro de atención especializado frente a uno generalista, un hospital frente a atención primaria, un gestor personal frente a una aplicación. Antes de redistribuir, quieren saber cuánto se perdería al mover a un grupo hacia el canal básico. Un experimento aleatorizado lo respondería, pero suele ser caro, lento o inviable; lo disponible son datos observacionales de cada canal. El enfoque estándar repondera a las unidades que hoy atiende el canal básico para que se parezcan al grupo objetivo, con pesos estimados por un clasificador. Formalmente, es un cambio de medida.

Este artículo trata de un rasgo de los datos de canal que el enfoque estándar ignora: \emph{cada canal registra las covariables a su manera}. Un gestor personal apunta un historial comercial que la aplicación nunca guarda; un hospital hace pruebas que atención primaria no hace; un canal registra lo que el cliente declara y otro lo mide. Mostramos que esto rompe la reponderación de dos formas, una visible y otra invisible.
\begin{itemize}
  \item \textbf{Una variable que solo registra un canal} rompe el soporte por construcción, por comparables que sean las unidades: la masa sin soporte es al menos la fracción del grupo objetivo con dato, y el clasificador ideal lo delata al estar seguro ($p=1$) en esos perfiles (Proposición~\ref{prop:rompe}); en una referencia experimental (NSW frente a CPS), una variable de ingresos censurada rompe el soporte en su extremo y los clasificadores estimados no lo detectan (sección~\ref{sec:lalonde}). Descartar la variable restaura la identificación cuando no predice el resultado en el canal básico dadas las demás covariables (Proposición~\ref{prop:descartar}); si lo predice, ni conservarla ni descartarla identifica el estimando.
  \item \textbf{Una variable medida con un instrumento distinto en cada canal} puede conservar el soporte. El estimador converge entonces a un valor desplazado, que derivamos exactamente, y que ningún análisis de los datos observados distingue de la verdad (Teorema~\ref{teo:medicion}). Con datos públicos en los que existen los dos instrumentos para las mismas personas (índice de masa corporal declarado y medido en NHANES), el desplazamiento derivado coincide de forma aproximada con el sesgo observado en 7 de 8 combinaciones de diseño y resultado (la excepción, la tensión con LightGBM, se explica por el modelo de resultado), y los diagnósticos habituales de solapamiento mejoran, respecto de la medición correcta, en lugar de delatarlo (sección~\ref{sec:nhanes}).
\end{itemize}

Para enunciar estos resultados hace falta la maquinaria de la reponderación en lenguaje de teoría de la medida (sección~\ref{sec:identificacion}). Casi toda es conocida y no la reclamamos: el estimando es la media contrafactual que entra en un efecto medio sobre los tratados, los pesos de odds $p/(1-p)$ son los de \citet{hirano2003}, estimar cocientes de densidades con clasificadores es habitual \citep{qin1998,bickel2009,sugiyama2012}, y la teoría doblemente robusta con \emph{cross-fitting} es la de \citet{chernozhukov2018}. Dos lecturas son nuevas, y las secciones~\ref{sec:canal} y~\ref{sec:instrumentos} se apoyan en ellas. La parte sin soporte de la ley objetivo es exactamente su restricción a $\{p=1\}$, que es la construcción de von Neumann de la descomposición de Lebesgue \citep[teorema 6.10]{rudin1987} leída en el clasificador; y sin soporte, el estimador doblemente robusto rellena esa parte en silencio extrapolando el modelo de resultado (Teorema~\ref{teo:lebesgue} y Corolario~\ref{cor:drsinsoporte}).

Damos además herramientas para la práctica. Cuando el grupo objetivo es un corte en una puntuación conocida, los pesos son la propensión de canal ordinaria con el corte impuesto (Proposición~\ref{prop:corte}), como en el análisis por subgrupos \citep{rassen2012}; estimar esa propensión con todas las unidades del canal premium redujo la dispersión hasta un 59\,\% en nuestra simulación cuando el solapamiento era bueno, pero no cuando era malo. Damos el sesgo exacto del recorte por tope de los pesos verdaderos con normalización (Proposición~\ref{prop:recorte}), que complementa a \citet{ma2020}, y adaptamos las cotas del modelo de sensibilidad marginal \citep{tan2006,dorn2023} a este estimando, con una calibración mediante el modelo de canal (sección~\ref{sec:sensibilidad}). Además de los dos análisis con datos públicos, una simulación contrasta los demás resultados: en cada escenario calculamos el límite de cada estimador, y registramos otras predicciones, antes de ejecutarlo, e informamos de las predicciones que fallaron (secciones~\ref{sec:diseno} y~\ref{sec:resultados}). El recorte con pesos estimados y el comportamiento del bootstrap son el objeto de una nota complementaria en preparación; aquí solo resumimos los resultados.

La sección~\ref{sec:planteamiento} plantea el problema y la sección~\ref{sec:identificacion}, la maquinaria de la reponderación. Las secciones~\ref{sec:canal} y~\ref{sec:instrumentos} contienen los dos fallos, cada uno seguido de datos públicos: el visible con la referencia NSW/CPS (sección~\ref{sec:lalonde}) y el invisible con NHANES (sección~\ref{sec:nhanes}). Las secciones~\ref{sec:corte}--\ref{sec:sensibilidad} reúnen herramientas para la práctica que no dependen del problema de medición: pesos con corte conocido, recorte e inferencia, y cotas de sensibilidad. Las secciones~\ref{sec:diseno} y~\ref{sec:resultados} presentan la simulación, y la sección~\ref{sec:discusion} concluye.

Usamos el lenguaje de la teoría de la medida (derivadas de Radon--Nikodym) en lugar de cocientes de densidades. No es una aportación, sino la formulación estándar; permite tratar covariables mixtas (continuas, categóricas y con valores ausentes) sin suponer que existen densidades.

\paragraph{Trabajo relacionado.} El mismo estimando, los pesos de odds y la misma función de influencia doblemente robusta aparecen en la literatura sobre generalizar y transportar resultados de ensayos a una población objetivo \citep{westreich2017,dahabreh2019}. \citet{heckman1997} encontraron empíricamente que el emparejamiento elimina gran parte del sesgo cuando el grupo de comparación procede del mismo mercado y responde el mismo cuestionario que los participantes; el Teorema~\ref{teo:medicion} formaliza esa lección como un resultado de indistinguibilidad. \citet{cook2008} y \citet{steiner2011} documentan también la importancia de medir las covariables igual y con fiabilidad, y que el error de medición diferencial no se corrige sin información externa es conocido en la literatura de error de medición; \citet{imai2010} estudian la identificación no paramétrica y el análisis de sensibilidad con error de medición diferencial del tratamiento. Lo más cercano al Teorema~\ref{teo:medicion} es la literatura sobre error de medición de covariables en métodos de puntuación de propensión. \citet{hong2017} tratan una covariable medida de forma distinta en los dos grupos de tratamiento y la corrigen con modelos bayesianos cuyas distribuciones a priori sobre el error actúan como parámetros de sensibilidad; \citet{webbvargas2017} imputan la covariable verdadera a partir de una muestra externa de calibración en la que se miden las dos versiones; y \citet{rudolph2016} adaptan análisis de sensibilidad para la confusión no observada al error de medición de covariables. Estos métodos corrigen el sesgo a partir de un modelo del error o de datos externos. Nuestros resultados son complementarios: caracterizamos el límite de los estimadores ponderado y doblemente robusto con instrumentos propios de cada canal, mostramos que sin esa información el sesgo no está identificado, que es por lo que esas correcciones necesitan distribuciones a priori o datos de validación, tratamos el caso límite de una variable que solo registra un canal, y mostramos que los diagnósticos habituales de solapamiento pueden mejorar a medida que crece el sesgo. Sobre covariables parcialmente observadas, el enfoque por patrón de ausencia \citep{dagostino2000,blake2020} exige que la variable no confunda allí donde falta; nuestras variables de canal son el caso en que el patrón lo fija el tratamiento, y la condición que usamos es un criterio de selección de covariables por el lado del resultado \citep{deluna2011,vanderweele2011}. Sobre diagnósticos, \citet{westreich2011} muestran que un AUC alto del modelo de propensión no es necesario ni suficiente para controlar la confusión; nuestra simulación añade que tampoco mide el solapamiento. Para diagnosticar fallos de positividad, véanse \citet{petersen2012} y \citet{damour2021}. Para el solapamiento limitado y el recorte, véanse \citet{crump2009}, \citet{ma2020,ma2023}, \citet{liu2023} y \citet{manamkoong2025}, cuya extrapolación de la parte sin solapamiento bajo un supuesto de suavidad complementa el Corolario~\ref{cor:drsinsoporte}; los comentamos donde aparecen (sección~\ref{sec:inferencia}).

\section{Planteamiento}\label{sec:planteamiento}

Trabajamos en un espacio de probabilidad $(\Omega,\mathcal F,\PP)$. Las covariables son una variable aleatoria $X\colon\Omega\to\X$, con $(\X,\A)$ un espacio estándar de Borel; esto garantiza distribuciones condicionadas regulares y admite vectores mixtos de variables continuas, categóricas y ausentes.

\begin{center}
\begin{tabular}{@{}ll@{}}
\toprule
Símbolo & Significado \\
\midrule
$D\in\{0,1\}$ & Canal que atiende a la unidad (1 = premium, 0 = básico) \\
$S\in\{0,1\}$ & Pertenencia al grupo objetivo, con $\{S=1\}\subseteq\{D=1\}$ \\
$Y(0),\,Y(1)$ & Resultados potenciales si la atiende el canal básico o el premium \\
$Y$ & Resultado observado \\
$P$ & Ley de $X$ en el grupo objetivo: $P(A)=\PP(X\in A\mid S=1)$ \\
$Q$ & Ley de $X$ en el canal donante: $Q(A)=\PP(X\in A\mid D=0)$ \\
\bottomrule
\end{tabular}
\end{center}

Suponemos $\PP(S=1)>0$ y $\PP(D=0)>0$. Un caso típico es $S=\ind\{D=1,\ s(X)\le c\}$: el grupo objetivo son las unidades del canal premium cuya puntuación $s(X)$ cae por debajo de un umbral, por ejemplo la cola de menor propensión de compra.

\paragraph{Estimando.} La media contrafactual del grupo objetivo en el canal básico,
\[
\theta=\E\bigl[Y(0)\mid S=1\bigr].
\]
La pérdida al desviar el grupo es $L=\E[Y\mid S=1]-\theta$; su primer término se observa directamente. Cuando las unidades generan varias observaciones, $P$, $Q$ y $\theta$ se toman a nivel de observación (sección~\ref{sec:inferencia}).

\paragraph{Regresiones.} Llamamos $m_0$ a una versión de $\E[Y\mid X,D=0]$, es decir, una función $\A$-medible y $Q$-integrable tal que, para todo $A\in\A$,
\[
\E\bigl[Y\,\ind\{X\in A\}\mid D=0\bigr]=\int_A m_0\dd Q .
\]
Existe por el teorema de Radon--Nikodym, aplicado a la medida $A\mapsto\E[Y\ind\{X\in A\}\mid D=0]$, que es absolutamente continua respecto de $Q$, y es única $Q$-casi seguro. Con $X$ discreta, $m_0(x)$ es el resultado medio del perfil $x$ (una tasa si $Y$ es binaria); la definición con conjuntos hace falta porque, con $X$ continua, cada valor concreto tiene probabilidad cero. Análogamente, $m_0^S$ es una versión de $\E[Y(0)\mid X,S=1]$, única $P$-casi seguro.

\paragraph{Supuestos.}
\begin{enumerate}[label=(A\arabic*)]
  \item \emph{Consistencia.} $Y=Y(0)$ sobre $\{D=0\}$. Falla si el canal básico cambia al recibir las unidades desviadas, por ejemplo por saturación.
  \item \emph{Intercambiabilidad condicional en media.} $m_0^S=m_0$ $P$-casi seguro, donde, por (A1), $m_0$ es también una versión de $\E[Y(0)\mid X,D=0]$. Solo concierne al canal básico: entre unidades con el mismo perfil, las del grupo objetivo obtendrían en él, en media, el mismo resultado que las unidades que el canal básico atiende hoy. No dice nada sobre $Y(1)$ y no es contrastable con los datos.
  \item \emph{Soporte.} $P\ll Q$: todo conjunto de perfiles con probabilidad cero en el canal donante tiene probabilidad cero en el grupo objetivo.
  \item \emph{Integrabilidad.} $\E[|Y(0)|\mid S=1]<\infty$ y $\E[|Y|\mid D=0]<\infty$; se cumple si $Y$ está acotada.
\end{enumerate}

\begin{observacion}[A2 presupone A3]\label{obs:a2a3}
Como $m_0$ solo está determinada $Q$-casi seguro, la igualdad $P$-casi seguro de (A2) depende de la versión elegida salvo que $P\ll Q$. Si $P$ tiene masa donde $Q$ no la tiene, $m_0$ es arbitraria allí y ningún dato del canal donante la fija. Este hecho reaparece en el Teorema~\ref{teo:lebesgue}: la parte del grupo objetivo sin soporte no es estimable, no solo difícil de estimar.
\end{observacion}

\begin{observacion}[Descomposición estructural de la pérdida]\label{obs:descomposicion}
Si el resultado es una suma de componentes, $Y=\sum_k Y_k$ (por ejemplo, compras de distintos productos), y el canal básico no puede producir las componentes de un conjunto $K_{\mathrm{ex}}$, es decir, $Y_k(0)=0$ casi seguro para $k\in K_{\mathrm{ex}}$, entonces, por linealidad,
\[
L=\sum_{k\in K_{\mathrm{ex}}}\E[Y_k\mid S=1]+\sum_{k\notin K_{\mathrm{ex}}}\bigl(\E[Y_k\mid S=1]-\E[Y_k(0)\mid S=1]\bigr).
\]
Leer $L$ como la pérdida causal exige consistencia también en el canal premium ($Y=Y(1)$ sobre $\{D=1\}$). El primer sumando se observa sin ningún supuesto de identificación; solo el segundo necesita las secciones~\ref{sec:identificacion} y~\ref{sec:clasificador}. Un error $\epsilon$ en $\theta_{\mathrm{com}}=\sum_{k\notin K_{\mathrm{ex}}}\E[Y_k(0)\mid S=1]$ produce un error $-\epsilon$ en $\hat L$, así que la conclusión sobre $L$ es poco sensible a fallos de (A2) o (A3) cuando $\theta_{\mathrm{com}}/L$ es pequeño. Una parte estructural grande no basta por sí sola: con $a_{\mathrm{ex}}$ y $a_{\mathrm{com}}$ las dos sumas de $\E[Y_k\mid S=1]$ y $\rho=a_{\mathrm{ex}}/L$, se tiene $\theta_{\mathrm{com}}/L=a_{\mathrm{com}}/L-(1-\rho)$, que puede ser grande aunque $\rho$ se acerque a 1. Dos advertencias: si la unidad obtiene después la componente por otra vía, $L$ es una cota superior de la pérdida para la organización; y con un resultado binario del tipo «compra algo» la descomposición no es aditiva, así que hay que trabajar con los indicadores por componente. El intercambio completo, que además inyecta en el canal premium un grupo procedente del básico, es la misma construcción con los papeles de los canales intercambiados.
\end{observacion}

\section{La maquinaria de la reponderación}\label{sec:identificacion}

Esta sección reúne las herramientas que usan las secciones~\ref{sec:canal} y~\ref{sec:instrumentos}: la identificación como cambio de medida, el clasificador que estima los pesos, la lectura del fallo de soporte en ese clasificador y los estimadores. Casi todo es conocido; el Teorema~\ref{teo:lebesgue} y el Corolario~\ref{cor:drsinsoporte} son las lecturas que aportamos.

\subsection{Identificación como cambio de medida}

Bajo (A1)--(A4), la media contrafactual $\theta$ es la media del resultado observado en el canal donante ponderada por $w=dP/dQ$. Ninguna parte del argumento usa densidades respecto de la medida de Lebesgue.

\begin{teorema}[Identificación]\label{teo:identificacion}
Supongamos (A1)--(A4) y además $\int_\X\E\bigl[|Y|\mid X=x,D=0\bigr]\,P(dx)<\infty$. Sea $w=dP/dQ$, que existe por (A3). Entonces
\[
\theta=\E\bigl[w(X)\,Y\mid D=0\bigr]=\frac{\E[\tilde w(X)\,Y\mid D=0]}{\E[\tilde w(X)\mid D=0]}\qquad\text{para todo }\tilde w=a\,w,\ a>0.
\]
\end{teorema}
La condición adicional se cumple siempre que $Y$ esté acotada.

\begin{proof}
\emph{Paso 1 (torre).} Como la ley de $X$ condicionada a $S=1$ es $P$, $\theta=\E[m_0^S(X)\mid S=1]=\int m_0^S\dd P$, finita por (A4), ya que $|m_0^S|\le\E[|Y(0)|\mid X,S=1]$ $P$-casi seguro.
\emph{Paso 2 (intercambiabilidad).} Por (A2), $\theta=\int m_0\dd P$.
\emph{Paso 3 (Radon--Nikodym).} Por (A3) existe $w\ge0$ medible con $P(A)=\int_A w\dd Q$ para todo $A$. La fórmula de cambio de medida, válida para indicadores, para funciones simples por linealidad y para funciones no negativas por convergencia monótona, se extiende a toda $g$ $P$-integrable: $\int g\dd P=\int g\,w\dd Q$. Con $g=m_0$, $\theta=\int m_0\,w\dd Q$.
\emph{Paso 4 (vuelta a los datos).} La condición adicional equivale a $\E[w(X)|Y|\mid D=0]<\infty$, así que $w(X)Y$ es integrable bajo $\PP(\cdot\mid D=0)$, y extrayendo lo conocido de la esperanza condicionada, $\E[w(X)Y\mid D=0]=\E[w(X)m_0(X)\mid D=0]=\int m_0\,w\dd Q=\theta$.
La forma autonormalizada se sigue de $\E[w(X)\mid D=0]=\int w\dd Q=P(\X)=1$: el cociente vale $\theta$ para $w$ y es invariante al multiplicar $w$ por $c>0$.
\end{proof}

Esta invarianza es la que permite, en la sección~\ref{sec:clasificador}, usar pesos conocidos solo salvo una constante.

\begin{ejemplo}\label{ej:juguete}
Sean tres perfiles A, B y C con $Q=(0,5;\,0,3;\,0,2)$, $P=(0,1;\,0,3;\,0,6)$ y $m_0=(4\,\%;\,2\,\%;\,0,5\,\%)$. La media ingenua del canal donante es $2,7\,\%$, pero $\theta=0,1\cdot4\,\%+0,3\cdot2\,\%+0,6\cdot0,5\,\%=1,3\,\%$. Los pesos $w=P/Q=(0,2;\,1;\,3)$ la recuperan desde el canal donante: $0,5\cdot0,2\cdot4\,\%+0,3\cdot1\cdot2\,\%+0,2\cdot3\cdot0,5\,\%=1,3\,\%$. La media ingenua duplica $\theta$ porque el canal donante tiene sobre todo perfiles A, que responden más.
\end{ejemplo}

\begin{proposicion}[Doble robustez]\label{prop:dr}
Bajo (A1)--(A4), sean $\bar w\ge0$ y $\bar m$ medibles, con las integrales siguientes finitas, y
\[
\psi(\bar w,\bar m)=\E\bigl[\bar w(X)\,(Y-\bar m(X))\mid D=0\bigr]+\int\bar m\dd P .
\]
Entonces $\psi(\bar w,\bar m)-\theta=\int_\X(\bar w-w)(m_0-\bar m)\dd Q$. En particular, $\psi(\bar w,\bar m)=\theta$ si $\bar w=w$ o $\bar m=m_0$, $Q$-casi seguro.
\end{proposicion}
\begin{proof}
Por el paso 4 de la prueba del Teorema~\ref{teo:identificacion} aplicado a $\bar w$, el primer sumando vale $\int\bar w(m_0-\bar m)\dd Q$. Por su paso 3, $\int\bar m\dd P-\theta=\int(\bar m-m_0)\,w\dd Q$. Sumando se obtiene el resultado.
\end{proof}

El error es un producto de los dos errores de estimación. Esa estructura es la que permite tasas lentas en cada modelo y, aun así, normalidad a ritmo $\sqrt n$ (Teorema~\ref{teo:normalidad}).

\subsection{El clasificador identifica la derivada de Radon--Nikodym}\label{sec:clasificador}

Un clasificador calibrado que distingue el grupo objetivo del canal donante da $w$ salvo una constante, y el conjunto donde ese clasificador está seguro, $\{p=1\}$, es exactamente la parte del grupo objetivo sin soporte.

\paragraph{Construcción.} Juntamos ambas poblaciones con una etiqueta: $(\tilde X,Z)$ con $Z\sim\mathrm{Bernoulli}(\pi)$, $\pi\in(0,1)$, $\tilde X\mid Z=1\sim P$ y $\tilde X\mid Z=0\sim Q$. La ley de $\tilde X$ es la mezcla $M=\pi P+(1-\pi)Q$, y el clasificador ideal es una versión $p(x)$ de $\PP(Z=1\mid\tilde X=x)$. En la práctica $\pi$ es la fracción del grupo objetivo en la muestra combinada; no afecta al estimando.

\begin{lema}\label{lem:mezcla}
$P\ll M$ y $Q\ll M$, con $dP/dM=p/\pi$ y $dQ/dM=(1-p)/(1-\pi)$, $M$-casi seguro.
\end{lema}
\begin{proof}
Si $M(A)=0$, entonces $\pi P(A)+(1-\pi)Q(A)=0$ y, como ambos coeficientes son positivos, $P(A)=Q(A)=0$. Además, para todo $A\in\A$, $\pi P(A)=\PP(Z=1,\tilde X\in A)=\int_A p\dd M$ por definición de esperanza condicionada; luego $p/\pi$ es una versión de $dP/dM$. El caso de $Q$ es idéntico con $1-p$.
\end{proof}

\begin{teorema}[Cociente de odds]\label{teo:odds}
Tomemos una versión de $p$ con valores en $[0,1]$ y sea $r=\frac{1-\pi}{\pi}\,\frac{p}{1-p}$ sobre $\{p<1\}$ y $r=0$ sobre $\{p=1\}$. Entonces $Q(\{p=1\})=0$ y, para todo $A\in\A$, $\int_A r\dd Q=P(A\cap\{p<1\})$. En particular, bajo (A3), $r=dP/dQ$ $Q$-casi seguro.
\end{teorema}
\begin{proof}
Por el Lema~\ref{lem:mezcla}, $Q(\{p=1\})=\int_{\{p=1\}}\frac{1-p}{1-\pi}\dd M=0$. Para $A\in\A$, por la regla de la cadena,
\[
\int_A r\dd Q=\int_{A\cap\{p<1\}}\frac{1-\pi}{\pi}\,\frac{p}{1-p}\cdot\frac{1-p}{1-\pi}\dd M=\int_{A\cap\{p<1\}}\frac{p}{\pi}\dd M=P(A\cap\{p<1\}).
\]
Bajo (A3), $P(\{p=1\})=0$ porque $Q(\{p=1\})=0$, así que $\int_A r\dd Q=P(A)$ para todo $A$.
\end{proof}

Combinado con la invarianza del Teorema~\ref{teo:identificacion}, bastan los odds $p/(1-p)$: la constante $(1-\pi)/\pi$ se cancela. Con un clasificador de log-odds $F$, el peso es $e^{F}$. El teorema exige que $p$ sea la probabilidad condicional verdadera: un clasificador bien ordenado pero mal calibrado ordena bien a las unidades y aun así da pesos sesgados.

\paragraph{Ejemplo~\ref{ej:juguete} (continuación).}
Con $\pi=0,2$, el clasificador ideal vale $p=(0,048;\,0,200;\,0,429)$ en A, B y C, con odds $(0,05;\,0,25;\,0,75)=w/4$, y el estimador autonormalizado da $1,3\,\%$. El clasificador $\tilde p=p^2$ ordena igual (mismo AUC), pero sus odds $(0,0023;\,0,042;\,0,225)$ dan $0,89\,\%$, un sesgo del 32\,\%. La calibración importa, no solo la discriminación.


\begin{teorema}[Soporte y descomposición de Lebesgue]\label{teo:lebesgue}
Sin suponer (A3):
\begin{enumerate}
  \item $P\ll Q$ si y solo si $P(\{p=1\})=0$.
  \item La descomposición de Lebesgue de $P$ respecto de $Q$ es $P=\Pac+\Ps$, con $\Pac(A)=P(A\cap\{p<1\})$, $\Ps(A)=P(A\cap\{p=1\})$ y $d\Pac/dQ=r$.
  \item Sea $\varepsilon=P(\{p=1\})$. Si $\varepsilon<1$, (A2) se cumple sobre $\{p<1\}$ y $\int_{\{p<1\}}\E[|Y|\mid X=x,D=0]\,P(dx)<\infty$, la ponderación autonormalizada con $r$ recupera
  \[
  \frac{\E[r(X)\,Y\mid D=0]}{\E[r(X)\mid D=0]}=\thac:=\E\bigl[Y(0)\mid S=1,\ p(X)<1\bigr],
  \]
  y $\theta=(1-\varepsilon)\thac+\varepsilon\,\ths$, con $\ths=\E[Y(0)\mid S=1,\ p(X)=1]$, definido cuando $\varepsilon>0$, no identificado.
\end{enumerate}
\end{teorema}
\begin{proof}
(2) Por el Teorema~\ref{teo:odds}, $\Pac$ tiene densidad $r$ respecto de $Q$, luego $\Pac\ll Q$; $\Ps$ está concentrada en $\{p=1\}$, que es $Q$-nulo, luego $\Ps\perp Q$; por la unicidad de la descomposición de Lebesgue, esta es la descomposición. El conjunto $\{p=1\}$ depende de la versión de $p$ solo módulo conjuntos $M$-nulos, que son nulos para $P$ y para $Q$. (1) $P\ll Q$ si y solo si $\Ps=0$. (3) Por el Teorema~\ref{teo:odds}, $\E[r(X)\mid D=0]=1-\varepsilon$; repitiendo los pasos 3 y 4 del Teorema~\ref{teo:identificacion} con $\Pac$, el numerador vale $\int_{\{p<1\}}m_0\dd P=(1-\varepsilon)\thac$. La relación con $\theta$ es la ley de probabilidad total.
\end{proof}

\begin{corolario}[Cotas de identificación parcial]\label{cor:cotas}
Con las hipótesis del Teorema~\ref{teo:lebesgue}(3), si $Y\in[0,1]$, entonces $(1-\varepsilon)\thac\le\theta\le(1-\varepsilon)\thac+\varepsilon$, y el sesgo de la ponderación es $\thac-\theta=\varepsilon(\thac-\ths)$.
\end{corolario}
La anchura del intervalo es $\varepsilon$, que puede superar al propio $\theta$ cuando el resultado es raro: si al Ejemplo~\ref{ej:juguete} se añade un perfil sin soporte con $P=(0,1;\,0,3;\,0,5;\,0,1)$, entonces $\varepsilon=0,1$, $\thac=1,39\,\%$ y el intervalo es $[1,25\,\%;\,11,25\,\%]$. Solo es informativo con un supuesto adicional; por ejemplo, si las unidades sin soporte no responden mejor que el resto del grupo, $\ths\le\thac$, el intervalo se reduce a $[(1-\varepsilon)\thac,\thac]$, aquí $[1,25\,\%;\,1,39\,\%]$, y el estimador ponderado es una cota superior de $\theta$.

\begin{corolario}[Estimador doblemente robusto sin soporte]\label{cor:drsinsoporte}
Bajo las hipótesis del Teorema~\ref{teo:lebesgue}, sea $\bar m$ $P$-integrable con $\bar m=m_0$ $P$-casi seguro sobre $\{p<1\}$. Con los pesos $r$ (el último término es $0$ cuando $\varepsilon=0$),
\[
\frac{\E[r(X)(Y-\bar m(X))\mid D=0]}{\E[r(X)\mid D=0]}+\int\bar m\dd P=(1-\varepsilon)\thac+\varepsilon\,\E_P\bigl[\bar m(X)\mid p(X)=1\bigr].
\]
\end{corolario}
\begin{proof}
Por el Teorema~\ref{teo:odds}, el numerador del primer sumando es $\int_{\{p<1\}}(m_0-\bar m)\dd P=0$. El segundo se parte en $\{p<1\}$, donde integra $(1-\varepsilon)\thac$, y $\{p=1\}$, donde integra $\varepsilon\,\E_P[\bar m\mid p=1]$.
\end{proof}

A diferencia del estimador ponderado, que se queda con la parte con soporte, el doblemente robusto rellena la parte sin soporte con la extrapolación del modelo de resultado. Recupera $\theta$ si y solo si esa extrapolación acierta en media, $\E_P[\bar m\mid p=1]=\ths$, que es un supuesto adicional no contrastable. \citet{manamkoong2025} proponen hacerla explícita bajo un supuesto de suavidad y acotar su efecto con un análisis de sensibilidad.

\paragraph{Lectura práctica.} En muestra finita nunca se observa $\hat p=1$ exactamente. Un diagnóstico natural es la masa estimada $\hat P(\hat p>1-t)$ para $t$ pequeño; un AUC cercano a 1 es una señal de alarma, porque indica que $P$ y $Q$ están cerca de ser mutuamente singulares, no que el clasificador sea bueno. Por debajo de ese extremo, el AUC no es una medida monótona del solapamiento: cuando el grupo objetivo es un corte en una puntuación, es alto incluso con buen solapamiento (sección~\ref{sec:diseno}).

\subsection{Estimadores}\label{sec:estimadores}

\paragraph{Muestreo.} Observamos una muestra combinada i.i.d. $O_i=(Z_i,X_i,(1-Z_i)Y_i)$, $i=1,\dots,n$, con $\PP(Z=1)=\pi$, $X\mid Z=1\sim P$ y $X\mid Z=0\sim Q$; para $\theta$ solo hace falta el resultado de las unidades donantes. Cuando una unidad (un cliente) genera varias observaciones (llamadas), $O_i$ es el bloque completo de esa unidad y todo lo que sigue se aplica a nivel de unidad. Suponemos que el canal y la pertenencia al grupo objetivo son propiedades de la unidad, constantes entre sus observaciones, y que $P$, $Q$, $m_0$ y $\theta$ se definen a nivel de observación, de modo que cada unidad pesa según su número de observaciones. Si ese número depende de $Y(0)$ dado $X$, los estimandos por unidad y por observación, y las versiones correspondientes de (A2), difieren.

\paragraph{Estimadores.} Con $\hat o=\hat p/(1-\hat p)$ y $n_1$ el número de observaciones del grupo objetivo,
\[
\hat\theta_{\mathrm{H}}=\frac{\sum_{i:Z_i=0}\hat o(X_i)\,Y_i}{\sum_{i:Z_i=0}\hat o(X_i)},\qquad
\hat\theta_{\mathrm{DR}}=\frac{\sum_{i:Z_i=0}\hat o(X_i)\,(Y_i-\hat m(X_i))}{\sum_{i:Z_i=0}\hat o(X_i)}+\frac1{n_1}\sum_{i:Z_i=1}\hat m(X_i).
\]
El primero es el estimador ponderado de Hájek. El segundo es el doblemente robusto (DR), la versión muestral de $\psi$ en la Proposición~\ref{prop:dr} con el término de pesos normalizado. Ambos usan \emph{cross-fitting}: las unidades se reparten en $K$ bloques y, para las observaciones del bloque $k$, $\hat p$ y $\hat m$ se entrenan con los demás.

\paragraph{Diagnósticos y evidencia.} De un clasificador ajustado se leen habitualmente tres diagnósticos: la masa sin soporte estimada $\hat\varepsilon=\hat P(\hat p>1-t)$ para un $t$ pequeño (Teorema~\ref{teo:lebesgue}); el AUC; y el tamaño muestral efectivo de Kish \citeyearpar{kish1965}, $\neff=(\sum_i\hat w_i)^2/\sum_i\hat w_i^2$ sobre las $n_0$ observaciones donantes, cuyo cociente $\neff/n_0$ estima $1/(1+\chi^2(P\|Q))$ (sección~\ref{sec:inferencia}). Las secciones~\ref{sec:canal} y~\ref{sec:instrumentos} estudian qué muestran estos diagnósticos cuando los canales registran las covariables de forma distinta. En todo el texto, $z$ denota (media $-$ predicción)/(error de Monte Carlo) entre réplicas, y una predicción está \emph{registrada} si se subió, con fecha, al repositorio público (sección~\ref{sec:diseno}) antes de la ejecución correspondiente.

\section{Variables que solo registra un canal: el fallo visible}\label{sec:canal}

Una covariable que solo se registra en uno de los canales rompe el soporte aunque las unidades sean comparables. Si no predice el resultado dadas las demás, descartarla restaura la identificación; si confunde, ni usarla ni descartarla identifica $\theta$.

\paragraph{Modelo de medición.} Cada unidad tiene covariables latentes $X^*=(V,C^*)$, donde $C^*$ toma valores en $\mathcal C\cup\{\vacio\}$ y $C^*=\vacio$ significa «sin dato real». Lo registrado depende del canal: $X=(V,C)$, con $C=C^*$ si $D=1$ y $C=\vacio$ si $D=0$. La sección~\ref{sec:instrumentos} trata el caso general $X=\varphi_D(X^*)$. Ejemplos: un historial comercial que solo guardan los gestores del canal premium, o pruebas diagnósticas que solo se hacen en el hospital. Escribimos $P^*$ y $Q^*$ para las leyes de $X^*$ en el grupo objetivo y en el canal donante.

\begin{lema}[Masa singular]\label{lem:singular}
$\varepsilon=\Ps(\X)=\sup\{P(B):\ B\in\A,\ Q(B)=0\}$.
\end{lema}
\begin{proof}
Sea $N$ un conjunto $Q$-nulo donde se concentra $\Ps$. Para todo $B$ con $Q(B)=0$, $\Pac(B)=0$, porque $\Pac\ll Q$; luego $P(B)=\Ps(B)\le\Ps(\X)$. Además $P(N)=\Pac(N)+\Ps(N)=\Ps(\X)$, así que el supremo se alcanza en $N$.
\end{proof}

\begin{proposicion}[Una variable de canal rompe el soporte]\label{prop:rompe}
Bajo el modelo de medición, sea $\kappa=\PP(C^*\ne\vacio\mid S=1)$ la fracción del grupo objetivo con dato real. Entonces $\varepsilon\ge\kappa$, y el clasificador ideal vale 1 $M$-casi seguro en los perfiles con $C\ne\vacio$. Si $\kappa>0$, (A3) falla con independencia de cómo sean las unidades.
\end{proposicion}
\begin{proof}
Sea $B=\mathcal V\times\mathcal C$, el conjunto de perfiles con dato registrado. Por el modelo de medición, $Q(B)=\PP(C\ne\vacio\mid D=0)=0$, mientras que $P(B)=\kappa$; el Lema~\ref{lem:singular} da $\varepsilon\ge\kappa$. Como $\PP(Z=0,\tilde X\in B)=(1-\pi)Q(B)=0$, $p=1$ $M$-casi seguro sobre $B$.
\end{proof}

En adelante, $g$ denota una versión de $\E[Y(0)\mid V,D=0]$, y $P_V$ y $Q_V$ las leyes de $V$ en el grupo objetivo y en el canal donante.

\begin{proposicion}[Descartar una variable que no predice el resultado]\label{prop:descartar}
Supongamos:
\begin{enumerate}
  \item \emph{Soporte latente:} $P^*\ll Q^*$.
  \item \emph{(A2) latente:} $\E[Y(0)\mid X^*,S=1]=\E[Y(0)\mid X^*,D=0]$, $P^*$-casi seguro.
  \item \emph{$C^*$ no predice el resultado en el canal básico dado $V$:} $\E[Y(0)\mid X^*,D=0]=\E[Y(0)\mid V,D=0]$, $Q^*$-casi seguro.
\end{enumerate}
Entonces, con $V$ como únicas covariables, (a) $P_V\ll Q_V$ y (b) $\E[Y(0)\mid V,S=1]=\E[Y(0)\mid V,D=0]$ $P_V$-casi seguro, así que los Teoremas~\ref{teo:identificacion}--\ref{teo:lebesgue} se aplican con $X=V$.
\end{proposicion}
\begin{proof}
(a) Si $Q_V(A)=0$, entonces $Q^*(A\times(\mathcal C\cup\{\vacio\}))=0$, y por (1) también $P^*(A\times(\mathcal C\cup\{\vacio\}))=P_V(A)=0$.
(b) Por (3), $\E[Y(0)\mid X^*,D=0]=g(V)$ salvo en un conjunto $Q^*$-nulo, que por (1) es también $P^*$-nulo. Por (2), $\E[Y(0)\mid X^*,S=1]=g(V)$ $P^*$-casi seguro, y como $V$ es función de $X^*$, la torre dentro de $\{S=1\}$ da $\E[Y(0)\mid V,S=1]=g(V)$.
\end{proof}

La condición (3) es un criterio de selección de covariables por el lado del resultado \citep{deluna2011,vanderweele2011}: basta aunque $C^*$ esté asociada al canal.

\paragraph{Cuando la variable sí confunde.} Si se mantienen (1) y (2) pero falla (3), sea $\mu^*(x^*)=\E[Y(0)\mid X^*=x^*,D=0]$; entonces $\theta=\int\mu^*\dd P^*$. Usando solo $V$, el estimador converge a $\int g\dd P_V=\int g(V)\dd P^*$, y su sesgo es
\[
\int g\dd P_V-\theta=-\int_{\X^*}(\mu^*-g)\dd P^*.
\]
Bajo $Q^*$ esa integral es cero por definición de $g$. Bajo $P^*$ solo puede ser distinta de cero si $C^*$ predice $Y(0)$ dado $V$ y además su distribución dado $V$ difiere entre grupos: son las dos condiciones de la definición clásica de confusor, necesarias pero no suficientes, porque los efectos podrían compensarse. La condición exacta para (A2) con $V$ como covariables es $\E_{P^*}[\mu^*-g\mid V]=0$ $P_V$-casi seguro; el sesgo en sí se anula si y solo si $\int(\mu^*-g)\dd P^*=0$, lo que también puede ocurrir por compensación. (3) es suficiente para ambas y no depende de $P^*$. Cuando el sesgo no es nulo, ninguna de las dos opciones identifica $\theta$: sin $C$ falla (A2), y con $C$ falla (A3) y, como se ve a continuación, en general también (A2) en la parte con soporte. Por eso tampoco sirven las cotas del Corolario~\ref{cor:cotas}, que suponen (A2) en esa parte. Lo que queda es un análisis de sensibilidad sobre (3).

\paragraph{Conservar la variable tampoco arregla la parte con soporte.} Con $C$ entre las covariables, la parte con soporte es $\{C=\vacio\}$. Pero «$\vacio$» no significa lo mismo en los dos grupos: en el canal donante mezcla unidades con $C^*=\vacio$ y con $C^*\ne\vacio$, porque ahí nunca se registra; en el grupo objetivo equivale a $C^*=\vacio$. Como todo el canal donante tiene $C=\vacio$, $m_0(v,\vacio)=g(v)$, y el estimador ponderado converge a $\E_{P^*}[g(V)\mid C^*=\vacio]$, que en general no es ni $\theta$ ni $\thac$. En efecto, (A2) en $\{C=\vacio\}$ exige $\mu^*(V,\vacio)=g(V)$ $P^*$-casi seguro sobre $\{C^*=\vacio\}$, es decir, que entre los donantes con el mismo $V$ respondan igual en el canal básico los que tienen historial y los que no; (3) lo implica, pero es más débil. El doblemente robusto, además, extrapola el modelo de resultado a las unidades con dato (Corolario~\ref{cor:drsinsoporte}). Ese modelo se entrena en el canal donante, donde $C$ nunca varía; si no depende de $C$ fuera del soporte, como ocurre con árboles, que no pueden partir por una variable constante en el entrenamiento, extrapola $g(V)$, y el doblemente robusto converge a $\E_{P^*}[g(V)]$, como si $C$ no se usara.

\paragraph{Contrastabilidad.} La condición (3) concierne a $Y(0)$ en el canal básico, donde $C^*$ no se observa, así que no es contrastable directamente. Que $C$ prediga el resultado en el canal premium es un indicio pero no una prueba, porque ahí se observa $Y(1)$.

\paragraph{Relación con el enfoque por patrón de ausencia.} \citet{dagostino2000} y \citet{blake2020} estiman la propensión dentro de cada patrón de valores ausentes y exigen que la variable parcialmente observada no confunda en el patrón donde falta. Una variable de canal es el caso en que el patrón lo determina el tratamiento: en el patrón observado no hay unidades donantes (Proposición~\ref{prop:rompe}), y en el ausente exige $\mu^*(V,\vacio)=g(V)$, que (3) implica.

\paragraph{Diagnóstico práctico.} Retirar una variable de canal debe reducir de forma visible el AUC del clasificador y la masa estimada $\hat P(\hat p>1-t)$; un salto grande en esa masa al incluir una variable es la firma de la Proposición~\ref{prop:rompe}. En la simulación, la masa sin soporte estimada con la variable de canal es $0,326$ frente a $\kappa=0,335$, y el estimador doblemente robusto se comporta como predice el Corolario~\ref{cor:drsinsoporte} y no como el ponderado (sección~\ref{sec:res-principal}).

\subsection{Ilustración: una referencia experimental con una covariable censurada}\label{sec:lalonde}

Una covariable no necesita faltar para romper el soporte: basta con que un canal la registre con un tope. Unos datos públicos con referencia experimental muestran un fallo así, y que el clasificador estimado no lo detecta.

\paragraph{Datos.} Usamos la submuestra de \citet{dehejia1999} del experimento National Supported Work, con 185 participantes y 260 controles aleatorizados, y el grupo de comparación CPS-1 de \citet{lalonde1986}, con 15\,992 personas. El grupo objetivo son los participantes, y los donantes, la muestra de la CPS. El estimando es la media de sus ingresos de 1978 sin el programa, cuya referencia experimental es la media de los controles: 4\,555 dólares (error típico 340). Las covariables son edad, años de estudio, raza, estado civil, falta de título, ingresos de 1974 y 1975 con indicadores de ingresos nulos, y los cuadrados de la edad y los años de estudio. Este análisis no se registró de antemano; damos la mediana sobre cinco particiones del \emph{cross-fitting}.

\paragraph{Resultados.} La media de la CPS sin ajustar es 14\,847 dólares. La ponderación sola da 5\,092 con un clasificador logístico y 3\,990 con LightGBM \citep{ke2017}, mientras que el estimador doblemente robusto da 4\,783 y 4\,716 (errores típicos 397 y 876), ambos a menos de un error típico de la referencia (340). El solapamiento es malo: con LightGBM, $\neff$ es 33 de 15\,992 donantes.

\paragraph{Tope en una de las fuentes.} Las dos fuentes registran los ingresos de forma distinta \citep{smith2005}. Una diferencia es visible: los ingresos de la CPS están censurados por arriba (el 19\,\% de los ingresos de 1974 está exactamente en el tope de 25\,862 dólares) y los del NSW no. Dos participantes (1,1\,\%) ganaron más que el tope en 1974, así que sus perfiles no tienen soporte por construcción: una masa singular en el sentido del Lema~\ref{lem:singular}, donde el clasificador ideal vale 1. Ninguno de los dos clasificadores estimados se acerca a 1: ambos dan a esas dos unidades un $\hat p$ inferior a 0,3, porque el 19\,\% de donantes en el tope hace que los ingresos altos parezcan propios de los donantes. Una falta de soporte en el extremo de una variable censurada apunta, así, el diagnóstico en la dirección equivocada. Aplicar el tope a los dos grupos cambia las estimaciones menos de 30 dólares, así que aquí no importa. El resultado también está censurado: el 26\,\% de los ingresos de 1978 de la CPS está en el tope, lo que sesga a la baja el modelo de resultado de los donantes; aplicar el tope a los controles baja la referencia en 54 dólares, a 4\,501, poco frente a los errores típicos. El tope es un instrumento propio de un canal que solo rompe el soporte en su extremo. La sección~\ref{sec:instrumentos} pasa a instrumentos que no lo rompen en absoluto; otras diferencias en cómo recogieron los ingresos las dos fuentes serían de ese tipo.

\section{Instrumentos de medida distintos: el sesgo invisible}\label{sec:instrumentos}

Una variable ausente en un canal es el caso más visible de un problema general: cada canal mide con su propio instrumento. Con la notación de la sección~\ref{sec:canal}, sea $X^*$ la covariable latente, con valores en un espacio estándar de Borel $\X^*$, y sean $P^*$ y $Q^*$ sus leyes en el grupo objetivo y en el canal donante. Cada canal registra $X=\varphi_D(X^*)$, con $\varphi_0,\varphi_1\colon\X^*\to\X$ medibles, de modo que $P=P^*\circ\varphi_1^{-1}$ y $Q=Q^*\circ\varphi_0^{-1}$. Sea $\mu^*(x^*)=\E[Y(0)\mid X^*=x^*,D=0]$.

\begin{teorema}[Instrumentos de medida distintos]\label{teo:medicion}
Supongamos (A1), soporte latente $P^*\ll Q^*$ y (A2) latente: $\E[Y(0)\mid X^*,S=1]=\mu^*(X^*)$ $P^*$-casi seguro.
\begin{enumerate}
  \item \emph{Variables comunes.} Para toda aplicación medible $\xi$ sobre $\X^*$, $P^*\circ\xi^{-1}\ll Q^*\circ\xi^{-1}$. En particular, si una $\upsilon$ medible sobre $\X$ cumple $\upsilon\circ\varphi_0=\upsilon\circ\varphi_1$ (variables que ambos canales miden igual), las leyes de $W=\upsilon(X)$ en los dos grupos cumplen $P_W\ll Q_W$.
  \item \emph{Identificación con las variables comunes.} Si además $\mu^*(X^*)=\E[Y(0)\mid W,D=0]$ $Q^*$-casi seguro, los Teoremas~\ref{teo:identificacion}--\ref{teo:lebesgue} se aplican con $W$ como covariables.
  \item \emph{Medición distinta con soporte.} Si $\varphi_0$ es inyectiva y $P\ll Q$ (en espacios estándar de Borel, $\varphi_0$ tiene entonces imagen boreliana e inversa medible por el teorema de Lusin--Souslin \citep[teorema 15.1]{kechris1995}, y $\varphi_1(X^*)$ está en esa imagen $P^*$-casi seguro porque $Q$ se concentra en ella), el estimador ponderado con $X$ converge, con $Y$ acotado, a
  \[
  \theta_X=\int_{\X^*}\mu^*\bigl(\varphi_0^{-1}(\varphi_1(x^*))\bigr)\,P^*(dx^*),\qquad
  \theta_X-\theta=\int_{\X^*}\bigl[\mu^*\circ\varphi_0^{-1}\circ\varphi_1-\mu^*\bigr]\dd P^*.
  \]
  \item \emph{Indistinguibilidad.} En las condiciones de 3, existe un mundo sin error de medición, con $\varphi_1=\varphi_0$ y la ley latente del grupo objetivo igual a $P^*\circ(\varphi_0^{-1}\circ\varphi_1)^{-1}$, que reproduce la ley de los datos del grupo objetivo y del canal donante (y la de todo lo observado si, además, las covariables registradas de todas las unidades premium tienen una ley absolutamente continua respecto de $Q$), cumple los supuestos mantenidos y en el que $\theta=\theta_X$. Ningún procedimiento basado solo en los datos observados distingue $\theta$ de $\theta_X$.
\end{enumerate}
\end{teorema}
\begin{proof}
(1) Si $Q^*(\xi^{-1}(B))=0$, por el soporte latente $P^*(\xi^{-1}(B))=0$. Si $\upsilon\circ\varphi_0=\upsilon\circ\varphi_1$, entonces $W=\upsilon(\varphi_D(X^*))=(\upsilon\circ\varphi_0)(X^*)$ en los dos canales, y basta aplicar lo anterior con $\xi=\upsilon\circ\varphi_0$.

(2) Es la Proposición~\ref{prop:descartar} con $V=W$, que es función de $X^*$.

(3) En el canal donante, $X=\varphi_0(X^*)$ y $\varphi_0$ es inyectiva con inversa medible, así que condicionar en $X$ equivale a condicionar en $X^*$; por (A1), $m_0(x)=\mu^*(\varphi_0^{-1}(x))$ $Q$-casi seguro, y como $P\ll Q$ esta versión queda fijada $P$-casi seguro. El estimador ponderado converge a $\int m_0\dd P$ (Teorema~\ref{teo:identificacion}, que para esta convergencia solo usa $P\ll Q$), y por el cambio de variables de la medida imagen, $\int m_0\dd P=\int m_0(\varphi_1(x^*))\,P^*(dx^*)=\theta_X$. Por (A2) latente y la torre, $\theta=\int\mu^*\dd P^*$.

(4) En el mundo alternativo, el canal donante no cambia. El grupo objetivo se registra con $\varphi_0$ y su ley latente es $\tilde P^*=P^*\circ(\varphi_0^{-1}\circ\varphi_1)^{-1}$, así que su ley observada es $P^*\circ(\varphi_0\circ\varphi_0^{-1}\circ\varphi_1)^{-1}=P^*\circ\varphi_1^{-1}=P$, la misma que en el mundo real. Para su resultado observado basta tomar, dado $X$, la misma ley condicional que en el mundo real, y con (A2) latente su respuesta en el canal básico es $\mu^*$; su $\theta$ es $\int\mu^*\dd\tilde P^*=\theta_X$. El mundo alternativo cumple los supuestos: como $\varphi_0$ es inyectiva, $Q^*(A)=Q(\varphi_0(A))$ y $\tilde P^*(A)=P(\varphi_0(A))$, así que $P\ll Q$ da $\tilde P^*\ll Q^*$, y (A2) latente vale por construcción. Con la condición adicional, las unidades premium que no están en el grupo objetivo se tratan igual, transformando su ley latente por $\varphi_0^{-1}\circ\varphi_1$, que entonces está definida casi seguro; sin ella, $\varphi_1(X^*)$ puede caer fuera de la imagen de $\varphi_0$ para esas unidades, y un mundo con $\varphi_1=\varphi_0$ no puede reproducirlas.
\end{proof}

\begin{observacion}[Instrumentos aleatorios]\label{rem:stochastic}
Los puntos 3 y 4 usan el instrumento del grupo objetivo solo a través de la ley $P$ de sus covariables registradas. Por tanto, valen literalmente cuando ese instrumento es aleatorio, $X=\varphi_1(X^*,\xi)$ con $\xi$ una variable aleatoria auxiliar, como el error de una declaración: si $\varphi_0$ es inyectiva y $P\ll Q$, el estimador ponderado converge a $\theta_X=\E[\mu^*(\varphi_0^{-1}(X))\mid S=1]$, y el mundo en el que las covariables latentes del grupo objetivo se distribuyen como $\varphi_0^{-1}(X)$ dado $S=1$ reproduce los datos del grupo objetivo y del canal donante. No hace falta ninguna condición sobre cómo se relaciona $\xi$ con $Y(0)$, porque el límite depende del resultado solo a través de $m_0$, que fija el canal donante.
\end{observacion}

\paragraph{Lectura.} Una variable medida de forma distinta es más peligrosa que una ausente. La ausente rompe el soporte y normalmente deja huella: el clasificador ideal está seguro, $p=1$, y $\hat\varepsilon>0$ (la excepción, un tope en una de las fuentes, se comenta en la sección~\ref{sec:lalonde}). La medida de forma distinta puede conservar el soporte, y el clasificador interpreta la diferencia de medida como una diferencia real entre unidades. En un ejemplo con $P^*=N(1;\,1)$, $Q^*=N(0;\,1)$, $\mu^*(x)=\sigma(-3-0,8x)$, con $\sigma(t)=1/(1+e^{-t})$, y un canal premium que registra la variable desplazada en $\Delta$, un desplazamiento de $\Delta=\pm0,3$ desviaciones típicas produce sesgos de $-20,5\,\%$ y $+25,4\,\%$ sobre $\theta=2,92\,\%$. En una pequeña simulación de este ejemplo (detalles en el repositorio), el estimador ponderado con clasificador logístico coincide con $\theta_X$ ($z=+0,6$ en ambos casos). El AUC cambia (0,76 sin error; 0,82 y 0,69 con él), pero de una forma que también produciría una diferencia real de población. La única salida es información externa al diseño: unidades medidas con ambos instrumentos, por ejemplo clientes que cambiaron de canal, que permiten estimar $\varphi_0^{-1}\circ\varphi_1$, como en la calibración externa de \citet{webbvargas2017}.

\subsection{Evidencia con datos públicos: índice de masa corporal declarado frente a medido}\label{sec:nhanes}

El Teorema~\ref{teo:medicion} se puede comprobar directamente con datos en los que existen, para las mismas personas, los dos instrumentos.

\paragraph{Datos.} Usamos el ciclo 2017--2018 de la Encuesta Nacional de Salud y Nutrición de EE.\,UU. \citep{nchs2020}: los 4\,034 adultos de 20 años o más con talla y peso medidos y declarados, tensión arterial sistólica (media de hasta tres lecturas), perímetro de cintura, nivel educativo y ratio de renta sobre el umbral de pobreza. No usamos los pesos de la encuesta y tratamos la muestra como la población. El IMC declarado se calcula con el peso y la talla declarados y se redondea a 0,1, la precisión con la que se publica el IMC medido, para que sean comparables. El IMC declarado es de media 0,81 menor que el medido (desviación típica 1,99), y la diferencia crece con el IMC: de 0,00 en el cuartil más bajo del IMC medido a $-1,66$ en el más alto, en línea con la literatura sobre antropometría declarada \citep{connorgorber2007}. Usarlo solo en un canal es una medición distinta en el sentido del Teorema~\ref{teo:medicion}; que el error no sea clásico (crece con el IMC) es lo que hace que el sesgo sea apreciable.

\paragraph{Diseño.} Cada réplica asigna a cada persona al grupo objetivo (un canal premium sintético) con probabilidad $\sigma\bigl(-1+\beta_b z(\mathrm{IMC})+0,5\,z(\mathrm{edad})-0,4\,z(\mathrm{renta})\bigr)$, donde el IMC es el medido, $z$ estandariza y $\beta_b=0,8$; alrededor del 30\,\% acaba en el grupo objetivo. El canal es sintético, pero las covariables, los resultados y el error de medición son reales. Como la asignación depende solo de las covariables verdaderas, la intercambiabilidad latente y el soporte se cumplen por construcción, y $\theta$, el resultado medio del grupo objetivo, se observa en cada réplica. Comparamos cuatro especificaciones, todas con edad, sexo, raza, nivel educativo y renta: (A) IMC medido en los dos grupos; (B) IMC declarado en el grupo objetivo y medido en los donantes, el caso de medición por canal del Teorema~\ref{teo:medicion}; (C) sin IMC; (D) IMC declarado en los dos grupos. El estimador es el doblemente robusto con \emph{cross-fitting} en cinco pliegues, con un clasificador logístico penalizado y una regresión ridge, ambos aditivos en splines cúbicos; hay 200 réplicas. Las réplicas comparten la población y solo varía la asignación, así que evaluamos el sesgo, no la cobertura.

\paragraph{Predicciones registradas.} Antes de ejecutar calculamos el límite de la especificación B con el Teorema~\ref{teo:medicion}(3), con $\mu$, ajustada con toda la muestra, una estimación de $m_0(v,b)=\E[Y\mid V=v,\mathrm{IMC}=b,D=0]$, que aquí coincide con $\E[Y\mid V=v,\mathrm{IMC}=b]$ porque la asignación depende solo de $(V,\mathrm{IMC})$: promediando sobre las probabilidades de asignación $a_i$, el sesgo predicho es $\sum_ia_i[\mu(V_i,b_i^{\mathrm{decl}})-\mu(V_i,b_i)]/\sum_ia_i$, es decir, $-0,47$\,mmHg en la tensión y $-2,44$\,cm en la cintura. El valor declarado no es función del medido, así que este es el caso de instrumento aleatorio de la Observación~\ref{rem:stochastic}: como la asignación depende solo de $(V,\mathrm{IMC})$, se cumple la intercambiabilidad latente, $\theta=\int m_0(V,\mathrm{IMC})\dd P^*$, y el estimador converge a $\int m_0(V,\mathrm{IMC}^{\mathrm{decl}})\dd P^*$. El criterio registrado fue $|z|<3$, admitiendo un error de modelo de hasta el 25\,\% del sesgo predicho. Predijimos además ausencia de sesgo en A, ninguna masa de $\hat p$ cercana a 1 en A ni en B, una diferencia de AUC entre A y B inferior a 0,02 y, en C, un sesgo negativo y mayor en valor absoluto que el de B en la cintura.

\begin{table}[ht]
\centering
\caption{NHANES, diseño base (200 réplicas). Sesgo medio del estimador DR con su error de Monte Carlo entre paréntesis. En ninguna réplica hubo unidades objetivo con $\hat p>0,99$.}\label{tab:nhanes}
\begin{tabular}{@{}lcccc@{}}
\toprule
Especificación & Tensión (mmHg) & Cintura (cm) & AUC & $\neff/n_0$ \\
\midrule
A: medido en los dos & $-0,11$ (0,06) & $-0,18$ (0,02) & 0,743 & 0,41 \\
B: medición distinta & $-0,48$ (0,06) & $-2,56$ (0,02) & 0,715 & 0,50 \\
\quad predicción (Teorema~\ref{teo:medicion}) & $-0,47$ & $-2,44$ & & \\
C: sin IMC & $-1,78$ (0,05) & $-10,70$ (0,04) & 0,653 & 0,74 \\
D: declarado en los dos & $-0,47$ (0,06) & $-1,32$ (0,03) & 0,737 & 0,44 \\
\bottomrule
\end{tabular}
\end{table}

\paragraph{Resultados.} La tabla~\ref{tab:nhanes} confirma de forma aproximada la predicción del Teorema~\ref{teo:medicion}: $-0,48$ frente a $-0,47$\,mmHg ($z=-0,3$) y $-2,56$ frente a $-2,44$\,cm. La desviación en la cintura, del 5\,\%, está dentro de la tolerancia registrada, pero es estadísticamente distinta de cero ($z=-5,7$); descontando el pequeño sesgo de A, $-2,37$ frente a $-2,44$. Nada en los datos observados delata el problema. Ninguna unidad objetivo tiene $\hat p>0,99$, y $\neff/n_0$ \emph{mejora} de 0,41 a 0,50: el grupo objetivo, seleccionado por IMC alto, declara un IMC desplazado hacia abajo, hacia los donantes, así que los grupos se separan menos y los pesos parecen mejores justo cuando el estimador está sesgado (el ruido solo, en D, lleva $\neff/n_0$ únicamente a 0,44). El AUC baja de 0,743 a 0,715, pero la especificación A no está disponible en la práctica, así que esa bajada no se puede observar. Quitar el IMC es mucho peor ($-10,70$\,cm), porque el IMC determina la asignación. Fallan dos predicciones registradas: A tiene un sesgo de $-0,18$\,cm en la cintura ($z=-8,3$), que crece con la fuerza de la asignación y desaparece con LightGBM ($-0,01$; tabla~\ref{tab:nhanes-robust}), así que procede de los modelos de splines; y el AUC baja más que el umbral registrado de 0,02.

La especificación D, registrada como exploratoria, sugiere una lección práctica con límites: usar el IMC declarado en \emph{los dos} grupos reduce a la mitad el sesgo de la cintura respecto de B, pero no el de la tensión ($-0,47$ frente a $-0,48$). Medir los dos grupos con el mismo instrumento impreciso deja confusión residual, pero no evalúa el modelo de resultado en un argumento desplazado; cuál de los dos efectos domina depende del resultado. La tabla~\ref{tab:nhanes-robust} muestra que en la cintura D fue mejor que B en los cuatro diseños, entre un 24 y un 73\,\%, mientras que en la tensión fue peor en dos de ellos. No tenemos un resultado general sobre cuándo se da este orden.

\begin{table}[ht]
\centering
\caption{NHANES, robustez (200 réplicas cada uno). Sesgo medio del estimador DR en la cintura (cm) y, en las dos últimas columnas, en la tensión; entre paréntesis, la predicción registrada.}\label{tab:nhanes-robust}
\resizebox{\textwidth}{!}{%
\begin{tabular}{@{}lcccccc@{}}
\toprule
Diseño & A & B & C & D & B, tensión (mmHg) & D, tensión (mmHg) \\
\midrule
Asignación débil ($\beta_b=0,4$) & $-0,08$ & $-2,26$ ($-2,19$) & $-5,69$ & $-0,61$ & $-0,42$ ($-0,41$) & $-0,27$ \\
Base ($\beta_b=0,8$), splines & $-0,18$ & $-2,56$ ($-2,44$) & $-10,70$ & $-1,32$ & $-0,48$ ($-0,47$) & $-0,47$ \\
Asignación fuerte ($\beta_b=1,2$) & $-0,46$ & $-2,89$ ($-2,61$) & $-14,44$ & $-2,20$ & $-0,56$ ($-0,51$) & $-0,69$ \\
Base, LightGBM & $-0,01$ & $-2,42$ ($-2,44$) & $-10,38$ & $-1,23$ & $-0,20$ ($-0,47$) & $-0,45$ \\
\bottomrule
\end{tabular}}
\end{table}

\paragraph{Robustez.} Variamos la fuerza con la que el IMC determina la asignación y sustituimos los splines por LightGBM, registrando de nuevo las predicciones antes de ejecutar (tabla~\ref{tab:nhanes-robust}). Para cada diseño registramos cinco: B dentro de la tolerancia anterior; el sesgo de A en la cintura por debajo del 20\,\% del de B (un criterio relativo, elegido después de que la ejecución base mostrara un sesgo de modelo en A; con $|z|<3$, A también fallaría en los diseños débil y fuerte); ninguna unidad objetivo con $\hat p>0,99$ en A ni en B; C negativo y mayor que B en valor absoluto en la cintura; y D menor que B en valor absoluto en la cintura. De esas 15 predicciones se cumplen 14. En los cuatro diseños, ninguna unidad objetivo tiene $\hat p>0,99$ y $\neff/n_0$ es mayor en B que en A. El fallo es la tensión con LightGBM: $-0,20$ frente a $-0,47$ ($z=4,0$). La tensión depende poco del IMC y tiene mucho ruido, así que LightGBM aplana esa dependencia y solo transmite parte del error de medición, mientras que la predicción usó una $\mu$ ajustada con splines. Una comprobación no registrada lo respalda: con $\mu$ ajustada con LightGBM, el sesgo predicho es $-0,31$, y el $-0,20$ observado queda a $z=1,7$ (en la cintura, $-2,40$ en lugar de $-2,44$). Cuando el resultado depende poco de la covariable mal medida, la propia predicción depende del modelo de resultado.

\section{Grupo objetivo definido por un corte conocido}\label{sec:corte}

Las secciones~\ref{sec:canal} y~\ref{sec:instrumentos} trataban de qué significan las covariables registradas. Esta sección y las dos siguientes toman las covariables como dadas y tratan de lo bien que se estiman los pesos y de cómo expresar la incertidumbre; esta empieza por la forma más habitual de definir el grupo objetivo.

En el caso típico, el grupo objetivo es un corte en una puntuación que la organización conoce, $S=D\cdot\ind\{s(X)\le c\}$. Entonces no hace falta un clasificador que distinga el grupo objetivo del canal donante: basta la propensión de canal ordinaria, con el corte impuesto a mano.

\begin{proposicion}[Pesos con corte conocido]\label{prop:corte}
Sea $S=D\cdot\ind\{s(X)\le c\}$, con $s$ medible y conocida, y sea $e$ una versión de $\PP(D=1\mid X)$. Bajo (A3), existe $k>0$ tal que
\[
\frac{dP}{dQ}(x)=k\,\ind\{s(x)\le c\}\,\frac{e(x)}{1-e(x)}\qquad Q\text{-casi seguro.}
\]
\end{proposicion}
\begin{proof}
Sea $\mu$ la ley de $X$ en la población. Como $s(X)$ es función de $X$, $\PP(S=1\mid X)=\ind\{s(X)\le c\}\,e(X)$. Por la definición de probabilidad condicionada, $P(A)=\int_A\ind\{s\le c\}\,e\dd\mu/\PP(S=1)$ y $Q(A)=\int_A(1-e)\dd\mu/\PP(D=0)$ para todo $A\in\A$, así que $P$ y $Q$ son absolutamente continuas respecto de $\mu$, con densidades proporcionales a $\ind\{s\le c\}\,e$ y a $1-e$. El conjunto $\{e=1\}$ es $Q$-nulo, y fuera de él, bajo (A3), la regla de la cadena da $dP/dQ=(dP/d\mu)/(dQ/d\mu)=k\,\ind\{s\le c\}\,e/(1-e)$, con $k=\PP(D=0)/\PP(S=1)$.
\end{proof}

La identidad es elemental, pero cambia el problema de estimación. El clasificador del Teorema~\ref{teo:odds} tiene que aprender un salto nítido en $\{s=c\}$, con los pocos casos del grupo objetivo como positivos. La propensión de canal suele ser una función suave y se estima con todas las unidades del canal premium. La constante $k$ se cancela al normalizar, y el diagnóstico de soporte del Teorema~\ref{teo:lebesgue} se lee en $\hat e$: falta soporte donde $\hat e\to1$ dentro de $\{s\le c\}$. La misma identidad, aplicada a las covariables ampliadas con una $U$ no observada, da el Lema~\ref{lem:calibracion} de calibración de la sección~\ref{sec:sensibilidad}.

\paragraph{Relación con la literatura.} La identidad no es nueva en sustancia. En el análisis por subgrupos con puntuaciones de propensión es práctica habitual estimar la propensión en toda la cohorte y aplicarla dentro del subgrupo \citep{rassen2012}, y \citet{dong2020} estudian cuándo conviene ese modelo y cuándo uno ajustado en cada subgrupo, que equilibra mejor dentro del subgrupo a costa de más varianza. Lo que añadimos es la comparación con el clasificador «grupo objetivo frente a donantes», el habitual cuando el cambio de medida se plantea como estimación de un cociente de densidades \citep{bickel2009,sugiyama2012}, en el caso en que el grupo objetivo es un corte en una puntuación de las propias covariables. Con modelos flexibles, los dos enfoques son consistentes y las diferencias son de muestra finita, pero afectan a cosas que la literatura de subgrupos no examina: el signo del recorte, el tamaño muestral efectivo y la validez del bootstrap (sección~\ref{sec:res-corte}).

\paragraph{Estimación.} En los estimadores de la sección~\ref{sec:estimadores}, $\hat o$ se sustituye por $\ind\{s\le c\}\,\hat e/(1-\hat e)$, y la muestra incluye también las unidades premium de fuera del grupo objetivo. Con muestreo poblacional de $(D,S,X,Y)$, la función de influencia correspondiente es $\PP(S=1)^{-1}\bigl[S(m_0-\theta)+(1-D)\,\ind\{s\le c\}\,e/(1-e)\,(Y-m_0)\bigr]$; esperamos un análogo del Teorema~\ref{teo:normalidad}, pero no lo demostramos.

\section{Recorte, varianza e inferencia}\label{sec:inferencia}

\subsection{Recorte de pesos}\label{sec:recorte}

Recortar los pesos cambia el estimando. Para $\tau>0$, sean $w_\tau=\min(w,\tau)$ y $\varrho_\tau(t)=\min(1,\tau/t)$ para $t>0$. Con los pesos verdaderos, el estimador autonormalizado recortado converge a
\[
\theta_\tau=\frac{\E[w_\tau(X)\,Y\mid D=0]}{\E[w_\tau(X)\mid D=0]}.
\]

\begin{proposicion}[Sesgo del recorte]\label{prop:recorte}
Bajo (A1)--(A4):
\begin{enumerate}
  \item $\displaystyle\theta_\tau-\theta=\frac{\Cov_P\bigl(m_0(X),\varrho_\tau(w(X))\bigr)}{\E_P[\varrho_\tau(w(X))]}$.
  \item Si $h(t)=\E_P[m_0(X)\mid w(X)=t]$ admite una versión no creciente, $\theta_\tau\ge\theta$; si admite una no decreciente, $\theta_\tau\le\theta$.
  \item Si $Y\ge0$, la versión sin normalizar está siempre por debajo: $\E[w_\tau(X)\,Y\mid D=0]\le\theta$.
\end{enumerate}
\end{proposicion}
\begin{proof}
Como $P(\{w=0\})=\int_{\{w=0\}}w\dd Q=0$, sobre $\{w>0\}$ se tiene $w_\tau=\varrho_\tau(w)\,w$. Por la torre y el cambio de medida, $\E[w_\tau(X)Y\mid D=0]=\int w_\tau m_0\dd Q=\E_P[\varrho_\tau(w)\,m_0]$ y $\E[w_\tau(X)\mid D=0]=\E_P[\varrho_\tau(w)]\in(0,1]$; restando $\theta=\E_P[m_0]$ se obtiene (1). Para (2), la torre da $\Cov_P(m_0,\varrho_\tau(w))=\Cov_P(h(T),\varrho_\tau(T))$ con $T=w(X)$. La función $\varrho_\tau$ es no creciente; si $h$ también lo es, $(h(T)-h(T'))(\varrho_\tau(T)-\varrho_\tau(T'))\ge0$ para una copia independiente $T'$, y su esperanza es $2\Cov_P(h(T),\varrho_\tau(T))$ (desigualdad de asociación de Chebyshev); el caso no decreciente es simétrico. Para (3), $\varrho_\tau\le1$ y $m_0\ge0$ dan $\E_P[\varrho_\tau(w)m_0]\le\theta$.
\end{proof}

\paragraph{Interpretación.} $w(x)$ grande significa que el perfil $x$ es mucho más frecuente en el grupo objetivo que en el canal donante. No hay un sentido por defecto: depende de si, dentro del grupo objetivo, los perfiles más típicos del canal premium responden más o menos en el canal básico. Si responden menos, $h$ es decreciente y el recorte sobreestima $\theta$ y, por tanto, infraestima la pérdida; en nuestra simulación responden más (sección~\ref{sec:res-recorte}). La condición trata de $h$ como función del peso, no de $X$, y es aproximadamente comprobable, porque $m_0$ está identificada en el soporte.

\paragraph{Ejemplo~\ref{ej:juguete} (recorte).} Con $\tau=2$, los pesos $(0,2;\,1;\,3)$ pasan a $(0,2;\,1;\,2)$ y $\theta_\tau=1,2\,\%/0,8=1,5\,\%$, frente a $\theta=1,3\,\%$. La fórmula da lo mismo: $\varrho_\tau=(1;\,1;\,2/3)$, $\E_P[\varrho_\tau]=0,8$ y la covarianza vale $0,16$ puntos, así que el sesgo es $0,2$ puntos. La versión sin normalizar da $1,2\,\%$, por debajo de $\theta$.

\begin{corolario}[Recorte con el estimador doblemente robusto]\label{cor:drrecorte}
Por la Proposición~\ref{prop:dr} con $\bar w=w_\tau$, $\psi(w_\tau,\bar m)-\theta=-\int_{\{w>\tau\}}(w-\tau)(m_0-\bar m)\dd Q$.
\end{corolario}
El corolario se refiere a la forma no normalizada $\psi$: el recorte solo sesga en la medida en que el modelo de resultado se equivoca en la región recortada. Para la forma normalizada que usamos, el sesgo es $\E_Q[w_\tau(m_0-\bar m)]/\E_Q[w_\tau]-\E_P[m_0-\bar m]$, que se anula si $\bar m=m_0$ en todo el soporte, pero no necesariamente si el modelo solo acierta en la región recortada. Como la Proposición~\ref{prop:dr} vale para cualesquiera pesos, al recortar pesos estimados el sesgo depende de los dos errores, y se anula si $\bar m=m_0$.

\paragraph{Pesos estimados.} En la práctica se recorta $\hat w$, y el recorte elimina a la vez peso verdadero y ruido de estimación, así que el signo de la Proposición~\ref{prop:recorte} no está garantizado. En una nota complementaria en preparación mostramos que se conserva si el error de estimación es un ruido multiplicativo independiente del perfil con densidad log-cóncava. Antes de recortar, el efecto del recorte calculado con $\hat m$ en lugar de $Y$, o $\hat m$ por deciles de masa de $\hat w$, indica en qué sentido empujará. \citet{crump2009} recortan por la propensión y redefinen el estimando; \citet{ma2020} mantienen el estimando y caracterizan el sesgo del recorte por descarte en el estimador sin normalizar, del que la Proposición~\ref{prop:recorte}(3) es la versión para el recorte por tope; la aportación de la Proposición~\ref{prop:recorte} es el caso normalizado, con su fórmula exacta. \citet{liu2023} tratan el truncamiento como un cambio de estimando para el efecto sobre los tratados.

\subsection{Varianza y normalidad asintótica}

\paragraph{Función de influencia.} Sean $o=p/(1-p)$ y $\sigma^2(x)=\Var(Y\mid X=x,D=0)$. La función de influencia del estimador doblemente robusto, con la misma forma que la eficiente para la media contrafactual de los tratados \citep{hahn1998}, es
\begin{gather*}
\phi(O)=\frac1\pi\Bigl[Z\,(m_0(X)-\theta)+(1-Z)\,o(X)\,(Y-m_0(X))\Bigr],\\
\mathcal V=\E[\phi^2]=\frac{\Var_P(m_0(X))}{\pi}+\frac{\E_Q[w(X)^2\sigma^2(X)]}{1-\pi},
\end{gather*}
donde el término cruzado se anula porque $Z(1-Z)=0$ y el segundo sumando usa $o=\pi w/(1-\pi)$ (Teorema~\ref{teo:odds}). Esta forma cerrada corresponde a observaciones i.i.d.; con varias observaciones por unidad, véase el apéndice.

\begin{proposicion}[Cuándo la varianza es finita]\label{prop:varianza}
Si $Y$ está acotada y $0<\sigma_-^2\le\sigma^2\le\sigma_+^2$ $Q$-casi seguro, entonces $\mathcal V<\infty$ si y solo si $\E_Q[w^2]<\infty$, es decir, si y solo si $\chi^2(P\|Q)=\E_Q[w^2]-1<\infty$. La continuidad absoluta $P\ll Q$ no implica esa condición.
\end{proposicion}
\begin{proof}
El primer sumando de $\mathcal V$ es finito porque $m_0$ está acotada; el segundo está entre $\sigma_-^2\,\E_Q[w^2]/(1-\pi)$ y $\sigma_+^2\,\E_Q[w^2]/(1-\pi)$. Para ver que $P\ll Q$ no basta, sea $Q$ uniforme en $(0,1)$ y $w(x)=x^{-1/2}/2$, que integra 1; entonces $P\ll Q$ pero $\E_Q[w^2]=\int_0^1(4x)^{-1}dx=\infty$.
\end{proof}

Así, (A3) basta para identificar $\theta$, pero no para estimarla a ritmo $n^{-1/2}$: con la varianza condicional acotada inferiormente, hace falta además $\chi^2(P\|Q)<\infty$, y la teoría estándar exige algo más fuerte, pesos acotados.

\paragraph{Tamaño muestral efectivo.} El $\neff$ de Kish (sección~\ref{sec:estimadores}) cumple con los pesos verdaderos $\neff/n_0\to1/\E_Q[w^2]=1/(1+\chi^2(P\|Q))$, así que estima la divergencia $\chi^2$; con pesos estimados puede moverse en los dos sentidos (el ruido de $\hat w$ lo reduce y el encogimiento de un clasificador regularizado lo aumenta), y en nuestra simulación queda por debajo del teórico salvo con solapamiento bajo, donde el propio valor teórico es impreciso (sección~\ref{sec:diseno}). En muestra finita siempre es finito, incluso cuando $\chi^2=\infty$, y entonces tiende a cero muy despacio y con mucho ruido: en el contraejemplo de la Proposición~\ref{prop:varianza}, la mediana sobre réplicas pasa de 0,45 con $n_0=10^3$ a unos 0,22 con $n_0=10^7$, y con $n_0=10^3$ los percentiles 1 y 99 entre réplicas son 0,05 y 0,65. Un $\neff$ razonable no garantiza que $\chi^2$ sea finita; conviene examinar además la estabilidad de los cuantiles altos de $\hat w$.

\begin{teorema}[Normalidad asintótica del estimador doblemente robusto]\label{teo:normalidad}
Supongamos (A1)--(A4), $Y$ acotada y además:
\begin{enumerate}
  \item \emph{Solapamiento fuerte:} $p(X)\le1-\eta$ $Q$-casi seguro para algún $\eta>0$, es decir, $w$ acotada.
  \item \emph{Cross-fitting} por unidad en un número fijo $K$ de bloques de tamaño proporcional a $n$, con $\hat p\le1-\eta$ y $\hat m$ acotada por las mismas cotas que $Y$.
  \item \emph{Tasas:} $r_o=\max_k\|\hat o_k-o\|$ y $r_m=\max_k\|\hat m_k-m_0\|$ en $L^2(Q)$ tienden a 0 en probabilidad, y $r_o\,r_m=o_P(n^{-1/2})$.
\end{enumerate}
Entonces $\sqrt n(\hat\theta_{\mathrm{DR}}-\theta)\to N(0,\mathcal V)$ en distribución, y la media muestral de $\hat\phi(O_i)^2$ estima $\mathcal V$ de forma consistente.
\end{teorema}
La prueba, en el apéndice~\ref{app:prueba}, es una adaptación del Teorema 3.1 de \citet{chernozhukov2018} y de su ejemplo del efecto sobre los tratados: por la Proposición~\ref{prop:dr}, el sesgo es bilineal en los errores de los parámetros auxiliares (ortogonalidad de Neyman), el \emph{cross-fitting} controla el término de proceso empírico sin condiciones de complejidad, y el denominador se trata con el método delta.

\paragraph{Más allá del solapamiento fuerte.} La condición 1 es más fuerte que (A3) y que $\chi^2<\infty$. Si los pesos no están acotados pero $\chi^2<\infty$, $\mathcal V$ es finita, aunque la prueba estándar no cubre el caso. Si $\chi^2=\infty$ y la varianza condicional está acotada inferiormente, ningún estimador regular converge a ritmo $\sqrt n$: es el caso de identificación irregular de \citet{khan2010}. \citet{ma2020} muestran que entonces el límite puede ser no normal y proponen inferencia por submuestreo, y \citet{ma2023} dan inferencia adaptativa para estimadores doblemente robustos con recorte y corrección de sesgo. Nuestra simulación (sección~\ref{sec:diseno}) está en ese régimen: su índice de canal, lineal en covariables normales, no está acotado dentro de la región objetivo, así que los pesos no están acotados aunque $\chi^2$ sea finita, y la implementación recorta $\hat p$ a $[10^{-6},1-10^{-6}]$ en lugar de a $1-\eta$; sus coberturas son evidencia fuera de las condiciones del Teorema~\ref{teo:normalidad}.

\paragraph{Intervalos.} Usamos $\hat\theta_{\mathrm{DR}}\pm1,96\,(\hat{\mathcal V}/n)^{1/2}$. El Teorema~\ref{teo:normalidad} no cubre el recorte de pesos ni el bootstrap, que evaluamos en la simulación (secciones~\ref{sec:res-recorte} y~\ref{sec:res-bootstrap}). Si el grupo objetivo se define con un umbral estimado $\hat c$, su incertidumbre añade un término a $\mathcal V$ que no incluimos: tratamos la regla de selección como fija.

\section{Sensibilidad a la confusión no observada}\label{sec:sensibilidad}

Si falla (A2), $\theta$ deja de estar identificado. El modelo de sensibilidad marginal acota cuánto puede alejarse $\theta$ del valor identificado en función de un único parámetro $\Lambda\ge1$.

\paragraph{Modelo.} Seguimos a \citet{tan2006} y \citet{dorn2023}; \citet{zhao2019} dan intervalos de sensibilidad para estimadores ponderados bajo este modelo mediante el bootstrap de percentiles. Supongamos que hay una variable no observada $U$ tal que la intercambiabilidad vale dado $(X,U)$: $\E[Y(0)\mid X,U,S=1]=\E[Y(0)\mid X,U,D=0]$. En la muestra combinada de la sección~\ref{sec:clasificador}, sea $\lambda(X,U)$ el cociente entre los odds de pertenecer al grupo objetivo dado $(X,U)$ y dado $X$ solo. El modelo MSM($\Lambda$) supone $\Lambda^{-1}\le\lambda(X,U)\le\Lambda$: conocer $U$ puede multiplicar o dividir esos odds como mucho por $\Lambda$. Con $\Lambda=1$ se recupera (A2).

\begin{proposicion}[Cotas de sensibilidad]\label{prop:sensibilidad}
Supongamos (A1), (A3), (A4), $Y\ge0$ y MSM($\Lambda$), y sea $\theta_0=\int m_0\dd P$ el valor identificado bajo (A2).
\begin{enumerate}
  \item \emph{Validez.} $\theta\in[\theta_0/\Lambda,\ \Lambda\theta_0]$.
  \item \emph{Exactitud para resultados raros.} Si además $\PP(Y>0\mid X,D=0)\le1/(1+\Lambda)$ $Q$-casi seguro, las dos cotas se alcanzan: existe un mundo compatible con MSM($\Lambda$) y con todo lo observado en el que $\theta=\Lambda\theta_0$, y otro en el que $\theta=\theta_0/\Lambda$.
  \item \emph{Caso binario general.} Para $Y\in\{0,1\}$ con $m=m_0(x)$, las cotas exactas integran respecto de $P$ los valores $\min\bigl(\Lambda m,\,1-(1-m)/\Lambda\bigr)$ (superior) y $\max\bigl(m/\Lambda,\,1-\Lambda(1-m)\bigr)$ (inferior); es la regla del cuantil de \citet{dorn2023} trasladada a $\theta$.
\end{enumerate}
\end{proposicion}
\begin{proof}
Por Bayes, $\lambda(x,\cdot)$ es la densidad de la ley de $U$ dado $X=x$ en el grupo objetivo respecto de la del canal donante, así que $\E_Q[\lambda\mid X]=1$, y la derivada de la ley de $(X,U)$ en el grupo objetivo respecto de la del donante es $w(X)\,\lambda(X,U)$. Aplicando el Teorema~\ref{teo:identificacion} a $(X,U)$, $\theta=\E_Q[w(X)\,\E_Q(\lambda Y\mid X)]$.

(1) Como $Y\ge0$ y $\Lambda^{-1}\le\lambda\le\Lambda$, $m_0(X)/\Lambda\le\E_Q[\lambda Y\mid X]\le\Lambda m_0(X)$; se integra contra $w\dd Q$.

(2) Sea $\pi_+(X)=\PP(Y>0\mid X,D=0)$. Para la cota superior, tómese $U=Y(0)$, $\lambda=\Lambda$ si $Y(0)>0$ y $\lambda=c(X)=(1-\Lambda\pi_+)/(1-\pi_+)$ si $Y(0)=0$. Entonces $\E_Q[\lambda\mid X]=1$; además $c\le1\le\Lambda$, y $c\ge1/\Lambda$ equivale a $\pi_+\le1/(1+\Lambda)$, así que MSM($\Lambda$) se cumple. Como $\E_Q[\lambda\mid X]=1$, ni la ley de $X$ en el grupo objetivo ni los datos del canal donante cambian, y $Y(0)$ del grupo objetivo nunca se observa, así que el mundo es compatible con lo observado; en él, $\E_Q[\lambda Y\mid X]=\Lambda m_0(X)$ y $\theta=\Lambda\theta_0$. Para la inferior, $\lambda=1/\Lambda$ si $Y(0)>0$ y $\lambda=(1-\pi_+/\Lambda)/(1-\pi_+)$ si $Y(0)=0$; este valor es $\ge1$ y es $\le\Lambda$ si $\pi_+\le\Lambda/(1+\Lambda)$, que se cumple porque $\pi_+\le1/(1+\Lambda)\le\Lambda/(1+\Lambda)$.

(3) Sean $\bar\lambda_1$ y $\bar\lambda_0$ las medias de $\lambda$ sobre $\{Y=1\}$ y $\{Y=0\}$ dado $X$. La restricción $\E_Q[\lambda\mid X]=1$ es $m\bar\lambda_1+(1-m)\bar\lambda_0=1$ con $\bar\lambda_0,\bar\lambda_1\in[1/\Lambda,\Lambda]$, y $\E_Q[\lambda Y\mid X]=m\bar\lambda_1$. Despejando $\bar\lambda_0$, el máximo de $m\bar\lambda_1$ es $\min(\Lambda m,\,1-(1-m)/\Lambda)$ y el mínimo, $\max(m/\Lambda,\,1-\Lambda(1-m))$; ambos se alcanzan con $\lambda$ constante en cada clase. Comprobamos además la fórmula frente a la solución por programación lineal en 2\,000 casos aleatorios, con error máximo de $10^{-15}$.
\end{proof}

\paragraph{Uso práctico.} Con un intervalo de confianza $[\hat L,\hat U]$ para $\theta_0$, el intervalo $[\hat L/\Lambda,\ \hat U\Lambda]$ combina la incertidumbre de muestreo con la de confusión. El \emph{$\Lambda$ de ruptura} es el menor $\Lambda$ que cambia la conclusión; por ejemplo, la pérdida $\E[Y\mid S=1]-\theta$ sigue siendo positiva para todo $\Lambda<\E[Y\mid S=1]/\theta_0$.

\paragraph{Calibración de $\Lambda$.} La comparación con covariables observadas, quitando cada una del modelo de propensión y midiendo cuánto cambian los odds, es una práctica establecida \citep{mcclean2024,baitairian2025}. Cuando el grupo objetivo es un corte en una puntuación de las propias covariables hay que hacerla con el modelo de canal, no con el de «grupo objetivo frente a donantes».

\begin{lema}[Calibración con el modelo de canal]\label{lem:calibracion}
Sea $S=D\cdot\ind\{s(X)\le c\}$ con $s$ observada. En la región $R=\{s(X)\le c\}$,
\[
\lambda(X,U)=\frac{\odds(D=1\mid X,U)}{\odds(D=1\mid X)}.
\]
\end{lema}
\begin{proof}
En la muestra combinada, por Bayes, $\odds(Z=1\mid x,u)=\frac{\pi}{1-\pi}\cdot\frac{\PP(S=1\mid x,u)/\PP(S=1)}{\PP(D=0\mid x,u)/\PP(D=0)}$, y lo mismo sin $u$. En $R$, $\PP(S=1\mid x,u)=\PP(D=1\mid x,u)$ y $\PP(S=1\mid x)=\PP(D=1\mid x)$; las constantes se cancelan en el cociente.
\end{proof}

Fuera de $R$, el modelo del grupo objetivo está dominado por el corte, que no es confusión. Dentro de $R$, el lema permite calibrar con un modelo de propensión de canal ordinario. Definimos el \emph{benchmark} de la covariable $j$ como $\Lambda_j(q)$, el cuantil $q$ en $R$ de $\max(\lambda_j,1/\lambda_j)$, con $\lambda_j=\odds(D=1\mid X)/\odds(D=1\mid X\text{ sin }j)$, ambos estimados fuera de muestra. Usamos cuantiles y no el máximo muestral, que crece con el tamaño de la muestra; son una escala de comparación, no una verificación del modelo. Dos advertencias: que un confusor oculto sea más débil que la covariable observada más fuerte es un supuesto, no una consecuencia de los datos; y quitar las covariables de una en una subestima la confusión si varias correlacionadas la reparten \citep{mcclean2024}.

\section{Diseño de la simulación}\label{sec:diseno}

La simulación usa un generador con $\theta$ conocida y cambia un factor cada vez sobre un escenario base, de modo que cada resultado teórico tiene un escenario que lo contrasta con una predicción numérica previa.

\paragraph{Generador.}
\begin{enumerate}
  \item \emph{Covariables:} $V=(V_1,\dots,V_4)$ normal con correlaciones $0,3^{|j-k|}$, más $V_5$ categórica de 3 niveles cuya distribución depende de $V_1$; todas son características del cliente. Una variable $U$ no observada, independiente de $V$, solo entra en el escenario de confusión no observada.
  \item \emph{Variable de canal latente:} $H=\ind\{C^*\ne\vacio\}$ indica historial real, con $\PP(H=1\mid V)=\sigma(-0,4+0,8V_1+0,3V_3)$; si $H=1$, el historial vale $N(0,5V_1;\,1)$. Se registra si $D=1$ y $H=1$, y $C=\vacio$ en otro caso.
  \item \emph{Canal:} $\PP(D=1\mid V,H,U)=\sigma(\alpha_0+\omega\,\alpha^\top V+\zeta H+\nu U)$. El parámetro $\omega$ controla el solapamiento, y $\zeta\ne0$ hace que la distribución de $C^*$ dado $V$ difiera entre canales, condición necesaria para que $C^*$ pueda confundir.
  \item \emph{Resultados:} dos productos comunes, $\PP(Y_k(d)=1\mid X^*,U)=\sigma(\beta_{k0}+\gamma_k d+\beta_k^\top V+\delta H+\nu'U+b)$ para $k=1,2$, con $\sigma$ la función logística y $b$ un efecto aleatorio de cliente, y un producto exclusivo del canal premium con $Y_3(0)=0$. El resultado que usan los estimadores es $Y=Y_1+Y_2$.
  \item \emph{Grupo objetivo:} $S=\ind\{D=1,\ s^*(V)\le c^*\}$, con $s^*$ el índice verdadero de propensión de compra en el canal premium y $c^*$ su percentil 10 en ese canal. La regla es fija, así que $\theta$ no cambia entre réplicas.
  \item \emph{Unidades y observaciones:} cada cliente genera entre 1 y 5 llamadas, en número independiente de todas las demás variables, con el canal y las covariables constantes entre ellas; por eso los estimandos por unidad y por observación coinciden.
  \item \emph{Verdad:} $\theta$ y los límites teóricos de cada estimador se calculan en una población de 1,5 millones de clientes; las medias condicionales verdaderas y $dP_V/dQ_V$ se obtienen integrando $H$, $U$ y el efecto de cliente por cuadratura de Gauss--Hermite. Su error de Monte Carlo, de unos 0,001 puntos (0,006 con confusión a través del historial y 0,015 con confusión no observada), es despreciable frente al de las medias de las réplicas. Los efectos del recorte, mucho más pequeños, se predicen como diferencias pareadas dentro de la población.
\end{enumerate}

\paragraph{Escenarios y predicciones.} El escenario base no incluye $C$, tiene $\zeta\ne0$ y $\delta=\nu=\nu'=0$, solapamiento alto ($\omega=0,5$; $\chi^2=3,6$), LightGBM \citep{ke2017} con \emph{cross-fitting} de 5 bloques por cliente para $\hat p$ y $\hat m$, y 30\,000 clientes por réplica: unas 62\,000 llamadas donantes y 2\,800 del grupo objetivo. Sobre él variamos la inclusión de $C$, la confusión por el historial ($\delta=1,2$), el solapamiento ($\omega=1,5$ y $3$; $\chi^2=6,3$ y unos 36; este último es inestable porque $w^2$ tiene una cola muy pesada: en seis poblaciones de 1,5 millones de clientes varió entre 33 y 63, así que el $\neff/n_0$ teórico con solapamiento bajo, 0,03, es solo orientativo), la especificación de los modelos, la confusión no observada ($\nu=\nu'=1$), el recorte en los percentiles 99, 95 y 90 de los pesos donantes y el tamaño (3\,000 a 30\,000 clientes). Cada escenario tiene 200 réplicas (las ejecuciones del bootstrap y del contraste de signo del recorte, menos; sección~\ref{sec:resultados}), con lo que la cobertura al 95\,\% se estima con un error típico de unos 1,5 puntos. Las predicciones que no se derivan directamente de la teoría se registraron con fecha antes de ejecutar el escenario correspondiente; el código, las predicciones con la fecha de su commit y los resultados brutos están disponibles en \url{https://github.com/pablotgp/channel-measurement}. Registrar significa aquí hacer un commit fechado en ese repositorio antes de la ejecución, no usar un registro externo.

\paragraph{El AUC no mide el solapamiento.} Al fijar $\omega$, el AUC del clasificador ideal se mantuvo entre 0,89 y 0,93 para $\omega$ entre 0,3 y 2, mientras $\chi^2$ variaba entre 3,5 y unos 11, sin crecer de forma monótona (es mínima para $\omega$ entre 0,5 y 1). La propia definición del grupo objetivo, un corte en una puntuación, ya lo separa del canal donante. Por eso fijamos los niveles por $\chi^2$, que es lo que gobierna la varianza (Proposición~\ref{prop:varianza}).

\section{Resultados de la simulación}\label{sec:resultados}

\subsection{Variables de canal, identificación y especificación}\label{sec:res-principal}

En los ocho escenarios de la tabla~\ref{tab:principal}, las medias de los estimadores quedan cerca de su límite predicho (todos los $|z|\le1,3$ salvo con solapamiento medio, que se comenta a continuación). Con muestras pequeñas fallan la dispersión y la cobertura, no la media (sección~\ref{sec:res-bootstrap}). Una discrepancia en el solapamiento medio ($z=-3,0$ en el Hájek y $-2,1$ en el DR) no reapareció en una réplica con 200 semillas nuevas ($z=+0,8$ y $+1,5$; con las 400 réplicas, $-1,4$ y $-0,3$). Esa réplica no estaba planificada de antemano; con unas 40 comparaciones, la probabilidad de algún $|z|\ge3$ por azar ronda el 10\,\%.

\begin{table}[ht]
\centering\small
\caption{Resultados principales. Tasas en \%; entre paréntesis, $z=$ (media $-$ predicción)/(error de Monte Carlo). DE: desviación típica del DR entre réplicas. Cobertura del intervalo de Wald del DR al 95\,\%, respecto de $\theta$. En los escenarios con $C$, el $\neff/n_0$ teórico se calcula con la parte con soporte; sale prácticamente igual que sin $C$ (0,217), porque $\PP(H=0\mid V,D=1)$ varía poco en la región objetivo.}\label{tab:principal}
\resizebox{\textwidth}{!}{%
\begin{tabular}{@{}lcccccc c@{}}
\toprule
Escenario & $\theta$ & Predicción & Hájek & DR & DE & Cobertura & $\neff/n_0$ (teórico) \\
\midrule
Base & 0,958 & 0,958 & 0,958 (0,0) & 0,961 (0,4) & 0,112 & 95,5\,\% & 0,14 (0,22) \\
Con $C$ & 0,958 & H 0,940; DR 0,958 & 0,939 ($-$0,1) & 0,958 (0,0) & 0,107 & 96,5\,\% & 0,12 (0,22) \\
$C$ confunde, sin $C$ & 1,712 & 1,314 & 1,313 ($-$0,1) & 1,320 (0,6) & 0,132 & 15,5\,\% & 0,14 (0,22) \\
$C$ confunde, con $C$ & 1,712 & H 1,281; DR 1,314 & 1,280 (0,0) & 1,320 (0,6) & 0,129 & 17,5\,\% & 0,12 (0,22) \\
Solapamiento medio & 1,178 & 1,178 & 1,151 ($-$3,0) & 1,158 ($-$2,1) & 0,136 & 94,5\,\% & 0,10 (0,14) \\
Solapamiento bajo & 1,274 & 1,274 & 1,283 (0,5) & 1,298 (1,1) & 0,300 & 91,0\,\% & 0,03 (0,03) \\
Confusión no observada & 2,697 & 1,193 & 1,198 (0,6) & 1,204 (1,3) & 0,127 & 0,0\,\% & 0,14 (0,20) \\
3\,000 clientes & 0,958 & 0,958 & 0,887 ($-$1,2) & 0,883 ($-$1,2) & 0,912 & 70,0\,\% & 0,02 (0,22) \\
\bottomrule
\end{tabular}}
\end{table}

\begin{itemize}
  \item \emph{Variable de canal.} Con $C$, la masa sin soporte estimada fue $\hat\varepsilon=0,326$, frente a $\kappa=0,335$ (Proposición~\ref{prop:rompe}), y el Hájek converge a $\thac$ (el límite de la ponderación de la sección~\ref{sec:canal} coincide aquí con $\thac$ porque $\delta=0$), con el sesgo predicho de $-0,018$ puntos. La predicción original asignaba también al DR el límite $\thac$ y falló ($z=2,4$ y $4,3$). El DR extrapola el modelo de resultado a la parte sin soporte (Corolario~\ref{cor:drsinsoporte}); la predicción derivada de ello (0,958 y 1,314) se añadió tras el escenario «Con $C$» y se registró antes de ejecutar «$C$ confunde, con $C$», donde se cumplió ($z=0,6$).
  \item \emph{Confusión por el historial.} Sin $C$, el sesgo es el predicho: $-0,40$ puntos, un 23\,\% de $\theta$. Con $C$ tampoco se recupera $\theta$, como dice la sección~\ref{sec:canal}.
  \item \emph{Identificación.} Sin confusión ni variable de canal, Hájek y DR recuperan $\theta$ y el intervalo cubre el 95,5\,\%.
  \item \emph{Confusión no observada.} Los estimadores convergen al límite predicho por los datos observables, a menos de la mitad de $\theta$, con cobertura nula. El error típico está bien calibrado (97\,\% de la dispersión real) y los diagnósticos se parecen a los del escenario base: nada en los datos delata el problema.
  \item \emph{Solapamiento.} Con solapamiento bajo ($\chi^2\approx36$, sección~\ref{sec:diseno}), la desviación típica casi se triplica, el error típico estimado se queda en el 87\,\% de la dispersión real y la cobertura baja al 91\,\%.
\end{itemize}

\begin{table}[ht]
\centering\small
\caption{Especificación (escenario base). Tasas en \%.}\label{tab:especificacion}
\begin{tabular}{@{}llllcc@{}}
\toprule
Clasificador & Modelo de resultado & Estimador & Predicción & Media ($z$) & Cobertura \\
\midrule
LightGBM & LightGBM & DR & 0,958 & 0,961 (0,4) & 95,5\,\% \\
LightGBM & lineal en $V_1$ & DR & 0,958 & 0,955 ($-$0,3) & 95,0\,\% \\
Logística lineal & LightGBM & DR & 0,958 & 0,946 ($-$2,3) & 88,0\,\% \\
Logística lineal & --- & Hájek & 1,024 & 1,023 ($-$0,4) & --- \\
Logística lineal & lineal en $V_1$ & DR & 1,157 & 1,150 ($-$1,5) & 13,5\,\% \\
--- & lineal en $V_1$ & Regresión & 1,069 & 1,075 (1,4) & --- \\
\bottomrule
\end{tabular}
\end{table}

Con el clasificador flexible y el modelo de resultado mal especificado (lineal), el DR no tiene sesgo y cubre el 95\,\%; con los dos mal, converge al límite predicho, lejos de $\theta$ (tabla~\ref{tab:especificacion}). Con el clasificador mal especificado (lineal) y el modelo de resultado flexible, queda un sesgo de $-0,012$ puntos y una cobertura del 88\,\%: el modelo de resultado es un LightGBM, no $m_0$, y el producto de errores no es despreciable a estos tamaños. Aun así, el DR reduce en un 82\,\% el sesgo del Hájek correspondiente.

\subsection{Recorte}\label{sec:res-recorte}

Con los pesos verdaderos y $m_0$, recortar en los percentiles 99, 95 y 90 baja la estimación en $0,37$, $3,13$ y $7,81$ milésimas de punto porcentual, como predice la Proposición~\ref{prop:recorte} ($0,37$, $3,18$ y $7,98$): los perfiles con más peso responden más. Con pesos estimados el signo se invierte ($+10,9\pm0,7$ en el percentil 90, positivo en las 60 réplicas). Dentro del DR, el efecto del recorte no se distingue de cero y reduce la raíz del error cuadrático medio (de 0,300 a 0,163 puntos con solapamiento bajo). Su error típico debe calcularse con los pesos sin recortar: con los recortados, la cobertura cae al 78,5\,\% con solapamiento bajo, y aun con los sin recortar es solo del 78--88,5\,\% con 3\,000 clientes. La inversión se analiza en la nota complementaria.

\subsection{Pesos con corte conocido}\label{sec:res-corte}

Comparamos el estimador con los pesos de la Proposición~\ref{prop:corte} con el estándar en las mismas réplicas: mismos datos, mismo LightGBM y mismo modelo de resultado; solo cambian los pesos. La propensión de canal se estima añadiendo las llamadas premium de fuera del grupo objetivo, que el estimador estándar no usa. Las predicciones se registraron antes de ejecutar estos escenarios.

\begin{table}[ht]
\centering\small
\caption{Pesos con corte conocido en el escenario base, por tamaño. 200 réplicas; $\theta=0,958\,\%$. DE en puntos. La última columna es el efecto del recorte del Hájek en el percentil 90 calculado con $m_0$, en milésimas, sobre estas 200 réplicas (el $+10,9$ del texto procede de otro conjunto de 60); con los pesos verdaderos es $-7,8$.}\label{tab:corte}
\begin{tabular}{@{}llccccc@{}}
\toprule
Clientes & Pesos & DR (\%) & DE & Cobertura & $\neff/n_0$ & Recorte p90 \\
\midrule
30\,000 & estándar & 0,961 & 0,112 & 95,5\,\% & 0,14 & $+10,4$ \\
30\,000 & corte conocido & 0,960 & 0,098 & 92,5\,\% & 0,20 & $-6,2$ \\
10\,000 & estándar & 0,962 & 0,294 & 88,0\,\% & 0,07 & $+19,8$ \\
10\,000 & corte conocido & 0,957 & 0,195 & 91,0\,\% & 0,17 & $-4,8$ \\
3\,000 & estándar & 0,883 & 0,912 & 70,0\,\% & 0,02 & $+50,0$ \\
3\,000 & corte conocido & 0,944 & 0,371 & 87,0\,\% & 0,11 & $-10,2$ \\
\bottomrule
\end{tabular}
\end{table}

Se cumplen cinco de las seis predicciones (tabla~\ref{tab:corte}): no hay sesgo apreciable; $\neff/n_0$ se acerca al 0,22 teórico; la dispersión baja un 12, un 34 y un 59\,\%; el recorte recupera el signo de los pesos verdaderos; y la cobertura mejora con 10\,000 y 3\,000 clientes. Falla la cobertura con 30\,000: predijimos entre el 93 y el 97\,\% y sale $92,5\,\%\pm1,9$, porque el error típico se queda en el 93\,\% de la dispersión real. Con 10\,000 clientes, la desviación típica queda un 15\,\% por encima de la que predice el ritmo $\sqrt n$ a partir del mismo estimador con 30\,000 ($0,098\cdot\sqrt3=0,170$), frente al 52\,\% del estándar.

Un análisis posterior, no registrado de antemano, separa las dos novedades: imponer el corte y usar las unidades premium de fuera del grupo objetivo. Imponer el corte sobre el clasificador estándar, sin datos nuevos, no reduce la dispersión (0,117, 0,298 y 0,933 puntos) y corrige solo en parte el signo del recorte ($+1,0$, $+2,7$ y $+9,6$). La mejora viene sobre todo de estimar la propensión de canal con muchos más casos positivos; el corte conocido es lo que permite usarla.

En los demás escenarios (apéndice~\ref{app:extra}), los estimadores coinciden con los mismos límites que el estándar, y la ganancia de precisión se mantiene con buen solapamiento pero desaparece con solapamiento medio y bajo.

\subsection{Sensibilidad}\label{sec:res-sensibilidad}

El $\Lambda$ mínimo para el que el intervalo poblacional $[\theta_0/\Lambda,\Lambda\theta_0]$ contiene $\theta$ es $\theta/\theta_0$: 1,30 cuando el historial confunde y 2,26 con confusión no observada. La cobertura del intervalo combinado $[\hat\ell/\Lambda,\hat u\Lambda]$ sube deprisa alrededor de ese valor (tabla~\ref{tab:sensibilidad}); queda por debajo de la aproximación normal porque el error típico es menor en las réplicas con estimación baja (correlación de 0,64 y 0,70). La condición de resultado raro de la Proposición~\ref{prop:sensibilidad}, comprobada mediante la cota $\PP(Y>0\mid V,D=0)\le g(V)=\E[Y\mid V,D=0]$ en el percentil 99,99 de $g$ en el canal donante (no en el supremo esencial), vale al menos hasta $\Lambda=5,6$ en el escenario base, 1,8 en el de confusión por el historial y 7,3 en el de confusión no observada. Con $U$ normal, $\lambda$ no está acotado y ningún $\Lambda$ finito cumple el modelo MSM en sentido estricto, así que en ese escenario la Proposición~\ref{prop:sensibilidad} no garantiza nada; $\Lambda=2,26$ contiene $\theta$ por construcción, y lo usamos solo como escala.

\begin{table}[ht]
\centering\small
\caption{Cobertura del intervalo combinado de sensibilidad.}\label{tab:sensibilidad}
\begin{tabular}{@{}lccccc@{}}
\toprule
Escenario & $\theta/\theta_0$ & $\Lambda=1,25$ & $\Lambda=1,5$ & $\Lambda=2$ & $\Lambda=2,5$ \\
\midrule
$C$ confunde, sin $C$ & 1,30 & 87\,\% & 100\,\% & 100\,\% & 100\,\% \\
Confusión no observada & 2,26 & 0\,\% & 3\,\% & 71\,\% & 99\,\% \\
\bottomrule
\end{tabular}
\end{table}

Para calibrar $\Lambda$, en el escenario de confusión no observada $U$ entra en el canal con coeficiente 1 y la covariable observada más fuerte, $V_1$, con 0,5. Predijimos, antes de calcularlo, que el método del Lema~\ref{lem:calibracion} recuperaría el valor exacto para $U$ y que $U$ superaría a todas las covariables observadas; ambas se cumplen (tabla~\ref{tab:calibracion}). El $\Lambda$ necesario para cubrir la verdad, 2,26 (conocido aquí solo porque conocemos $\theta$), queda entre la mediana y el cuantil 90 de $V_1$: basta un confusor algo más fuerte que la covariable observada más fuerte para explicar toda la diferencia. No es el $\Lambda$ de ruptura de la sección~\ref{sec:sensibilidad}, que depende de la conclusión de interés y se calcula con los datos. En una segunda estimación con otra muestra, el modelo «grupo objetivo frente a donantes» coincide con el de canal en la mediana (2,06 frente a 1,99), pero sobrestima los extremos: 6,34 y 29,98 en los cuantiles 90 y 99, frente a 5,37 y 15,51 con el de canal y 5,38 y 14,07 exactos.

\begin{table}[ht]
\centering\small
\caption{Calibración de $\Lambda$ con el modelo de canal: cuantiles en $R$ de $\max(\lambda,1/\lambda)$.}\label{tab:calibracion}
\begin{tabular}{@{}lccc@{}}
\toprule
Variable & Mediana & Cuantil 90 & Cuantil 99 \\
\midrule
$U$, exacto & 1,99 & 5,38 & 14,07 \\
$U$, estimado & 2,03 & 5,21 & 14,14 \\
$V_1$ & 1,72 & 3,07 & 5,81 \\
$V_2$ & 1,17 & 1,53 & 2,17 \\
$V_3$ & 1,14 & 1,45 & 1,98 \\
$V_4$ & 1,12 & 1,36 & 1,70 \\
$V_5$ & 1,08 & 1,21 & 1,42 \\
\bottomrule
\end{tabular}
\end{table}

\subsection{Tamaño muestral e intervalos: Wald y bootstrap}\label{sec:res-bootstrap}

Con el estimador estándar, la desviación típica del DR es 0,912, 0,508, 0,294, 0,183 y 0,112 puntos con 3\,000, 5\,000, 10\,000, 15\,000 y 30\,000 clientes: cae como $n^{-0,91}$, mucho más deprisa que $n^{-1/2}$, y la cobertura del intervalo de Wald es del 70, 80, 88, 92 y 95,5\,\%. Las predicciones de 5\,000 y 15\,000 se registraron de antemano y se cumplieron. El cociente $\neff/n_0$ estimado sigue la misma progresión (0,02, 0,07 y 0,14 con 3\,000, 10\,000 y 30\,000, frente al 0,22 teórico). Esto sugiere un diagnóstico, aún sin contrastar: si $\neff/n_0$ sigue creciendo al aumentar la muestra, los intervalos todavía no son fiables.

Evaluamos también el bootstrap por unidad (cliente), reestimando LightGBM y el \emph{cross-fitting} en cada remuestreo. Con los pesos estándar sobrestima la dispersión entre un 35 y un 60\,\% (en 80 réplicas, un error típico bootstrap medio de 0,165 frente a una dispersión real de 0,119 con 30\,000 clientes, y de 0,446 frente a 0,322 con 10\,000), y una variante bayesiana que pondera a los clientes en lugar de copiarlos solo reduce en parte el exceso. Con los pesos de corte conocido, mucho menos ruidosos, el bootstrap bayesiano queda casi calibrado: 0,101 frente a 0,096 y 0,208 frente a 0,187 en 40 réplicas, con coberturas del 95 y del 90\,\%. Fallaron las cuatro predicciones registradas sobre el bootstrap; la evidencia sugiere que el exceso está ligado al ruido de los pesos estimados, y dejamos los detalles para la nota complementaria.

\section{Discusión}\label{sec:discusion}

\paragraph{Recomendaciones prácticas.} De la teoría y de la simulación se siguen seis pautas para evaluar un cambio de canal con datos observacionales:
\begin{enumerate}
  \item Si el grupo objetivo es un corte en una puntuación conocida, estimar los pesos con la propensión de canal, usando todas las unidades del canal premium, e imponer el corte (Proposición~\ref{prop:corte}).
  \item Excluir de las covariables las variables que solo registra un canal cuando no predicen el resultado dadas las demás (Proposición~\ref{prop:descartar}), y diagnosticar la falta de soporte con la masa de $\hat p$ o $\hat e$ cercana a 1 (Proposición~\ref{prop:rompe}). Cuando la medición distinta conserva el soporte, ninguna comprobación con los datos la detecta (Teorema~\ref{teo:medicion}), y los diagnósticos de solapamiento pueden incluso mejorar (sección~\ref{sec:nhanes}); hace falta información externa. Si un canal mide una covariable con menos precisión y no se puede obtener la versión precisa, considerar usar también la menos precisa en los dos canales y comparar: en nuestra ilustración redujo el sesgo en la cintura pero no en la tensión, y no tenemos un resultado general.
  \item Usar el estimador doblemente robusto con \emph{cross-fitting} por unidad, sabiendo que sin soporte extrapola el modelo de resultado (Corolario~\ref{cor:drsinsoporte}).
  \item Para el intervalo, Wald con muestras grandes. Con muestras medianas y pesos de corte conocido, tanto Wald como el bootstrap bayesiano por unidad quedaron cerca del nivel nominal en la simulación (90--95\,\%, el bootstrap con solo 40 réplicas); con el clasificador estándar, el bootstrap es conservador. Desconfiar si $\neff/n_0$ sale muy por debajo de lo esperable.
  \item No trasladar el signo del recorte de los pesos verdaderos a los estimados sin comprobarlo. Si se recorta, hacerlo dentro del estimador doblemente robusto y calcular el error típico con los pesos sin recortar.
  \item Acompañar la estimación de un análisis de sensibilidad con $\Lambda$ calibrado frente a las covariables observadas mediante el modelo de canal (Lema~\ref{lem:calibracion}).
\end{enumerate}

\paragraph{Limitaciones.} La evidencia de la simulación procede de un único generador de datos, construido para que cada resultado teórico tuviera un escenario que lo contrastase, y las ilustraciones con datos públicos son semisintéticas (el canal lo asignamos nosotros) o se apoyan en una única referencia experimental; las conclusiones cuantitativas de la simulación, como las ganancias del corte conocido o el exceso del bootstrap, deben leerse como ilustraciones de los mecanismos, no como magnitudes generales. Las cotas de sensibilidad suponen el modelo MSM, que con confusores no acotados no se cumple en sentido estricto. Tratamos el umbral que define el grupo objetivo como fijo. Y el estimando es la media contrafactual con el canal básico tal como es hoy: si recibir las unidades desviadas lo satura, (A1) falla.

\paragraph{Problemas abiertos.} Cuatro preguntas quedan sin respuesta: una teoría que explique el exceso del bootstrap con pesos estimados ruidosos y la validez del bootstrap bayesiano con los pesos de corte conocido; la contribución a la varianza de un umbral estimado como cuantil de la puntuación, que se relaciona con la elección del propio umbral; la combinación del análisis de sensibilidad con la falta de soporte (Corolario~\ref{cor:cotas}; \citealp{manamkoong2025}) y su extensión a la condición (3) de la Proposición~\ref{prop:descartar}; y un resultado general sobre cuándo medir los dos grupos con el mismo instrumento impreciso es mejor que medirlos de forma distinta (sección~\ref{sec:nhanes}).


\appendix
\section{Prueba del Teorema~\ref{teo:normalidad}}\label{app:prueba}

El estimador doblemente robusto es $\theta$ más una media de términos i.i.d. de media cero, más restos de orden $o_P(n^{-1/2})$; el teorema central del límite hace el resto.

\paragraph{Notación.} Trabajamos primero con observaciones i.i.d. $O=(Z,X,(1-Z)Y)$; el caso con varias observaciones por unidad se trata al final. $P_nf=n^{-1}\sum_if(O_i)$; $\E f$ es la esperanza poblacional; $\|h\|_Q$ y $\|h\|_P$ son normas de $L^2(Q)$ y $L^2(P)$. Sean $o_{\max}=(1-\eta)/\eta$ y $w_{\max}=(1-\pi)o_{\max}/\pi$: por la condición 1 y el Teorema~\ref{teo:odds}, $o\le o_{\max}$ y $w\le w_{\max}$, y por la condición 2, $\hat o\le o_{\max}$. Las observaciones están acotadas: $|Y|\le B_Y$ y $|\hat m|\le B_Y$. En el bloque $k$, con $n_k$ observaciones, $(\hat o_k,\hat m_k)$ se entrenan con los datos de fuera del bloque; $\E_k$ y $\PP_k$ denotan esperanza y probabilidad condicionadas a esos datos, y $P_{n,k}$ la media empírica dentro del bloque.

\paragraph{Tres identidades.} Para toda $h$ medible y acotada:
\begin{align*}
&\text{(I1)}\quad \E[(1-Z)\,o(X)\,h(X)]=\pi\,\E_P[h],\\
&\text{(I2)}\quad \E[(1-Z)\,h(X)\,(Y-m_0(X))]=0,\\
&\text{(I3)}\quad \|h\|_P^2\le w_{\max}\,\|h\|_Q^2 .
\end{align*}
(I1): $\E[(1-Z)oh]=(1-\pi)\E_Q[oh]$, y $o=\pi w/(1-\pi)$ por el Teorema~\ref{teo:odds}, luego vale $\pi\E_Q[wh]=\pi\E_P[h]$. (I2): torre condicionando en $X$ sobre $\{Z=0\}$, donde $\E[Y\mid X,Z=0]=m_0(X)$. (I3): $\|h\|_P^2=\E_Q[wh^2]\le w_{\max}\,\E_Q[h^2]$.

\paragraph{Paso 1 (la función $\psi$ y su sesgo).} Para funciones acotadas $(o',m')$, sea $\psi(O;o',m')=(1-Z)\,o'(X)\,(Y-m'(X))+Z\,(m'(X)-\theta)$. Entonces $\pi\phi=\psi(\cdot;o,m_0)$ y
\[
\E[\psi(O;o',m')]=(1-\pi)\,\E_Q[(o'-o)(m_0-m')].
\]
Por (I2) con $h=o'$, $\E[(1-Z)o'(Y-m')]=(1-\pi)\E_Q[o'(m_0-m')]$. Por (A2), $\theta=\E_P[m_0]$, y por (I1) con $h=m'-m_0$, $\E[Z(m'-\theta)]=\pi\E_P[m'-m_0]=(1-\pi)\E_Q[o(m'-m_0)]$. Sumando se obtiene la fórmula; en particular, $\E\phi=0$. Es la Proposición~\ref{prop:dr} escrita con odds.

\paragraph{Paso 2 (linealización).} Con $\hat A=P_n[(1-Z)\hat o(Y-\hat m)]$, $\hat B=P_n[(1-Z)\hat o]$ y $\hat\pi=P_nZ$, el estimador es $\hat\theta=\hat A/\hat B+P_n[Z\hat m]/\hat\pi$. En el suceso $\{\hat\pi>\pi/2,\ \hat B>\pi/2\}$, cuya probabilidad tiende a 1 y en el que trabajamos en adelante, $P_n[Z\theta]/\hat\pi=\theta$ y el álgebra da exactamente
\[
\hat\theta-\theta=\frac1\pi P_n\psi(O;\hat o,\hat m)+\hat A\Bigl(\frac1{\hat B}-\frac1\pi\Bigr)+P_n[Z(\hat m-\theta)]\Bigl(\frac1{\hat\pi}-\frac1\pi\Bigr).
\]
Los pasos 3 y 4 muestran que el primer término es $P_n\phi+o_P(n^{-1/2})$; el paso 5, que los otros dos son $o_P(n^{-1/2})$.

\paragraph{Paso 3 (proceso empírico).} Sea $\Delta_k=\psi(\cdot;\hat o_k,\hat m_k)-\psi(\cdot;o,m_0)$. Condicionado a los datos de fuera del bloque $k$, $\hat o_k$ y $\hat m_k$ son funciones fijas y los $\Delta_k(O_i)$ del bloque son i.i.d. Por la desigualdad de Chebyshev condicional, para todo $\epsilon>0$,
\[
\PP_k\Bigl(\bigl|(P_{n,k}-\E_k)\Delta_k\bigr|>\epsilon\,n^{-1/2}\Bigr)\le\frac{n}{n_k}\,\frac{\E_k[\Delta_k^2]}{\epsilon^2}.
\]
Como $\Delta=(1-Z)[(\hat o-o)(Y-m_0)+\hat o(m_0-\hat m)]+Z(\hat m-m_0)$ y $|Y-m_0|\le2B_Y$, la desigualdad $(a+b+c)^2\le3(a^2+b^2+c^2)$ y (I3) dan
\[
\E_k[\Delta_k^2]\le3\Bigl[4B_Y^2(1-\pi)\,\|\hat o_k-o\|_Q^2+\bigl((1-\pi)o_{\max}^{\,2}+\pi w_{\max}\bigr)\|\hat m_k-m_0\|_Q^2\Bigr],
\]
que tiende a 0 en probabilidad por la condición 3. Como $n/n_k\to K$, la probabilidad condicional tiende a 0 en probabilidad, y al estar acotada por 1, la convergencia dominada da que la probabilidad sin condicionar también tiende a 0. Luego $(P_{n,k}-\E_k)\Delta_k=o_P(n^{-1/2})$ para cada uno de los $K$ bloques.

\paragraph{Paso 4 (sesgo).} Como $\E_k\psi(\cdot;o,m_0)=0$, el paso 1 y Cauchy--Schwarz dan
\[
|\E_k\Delta_k|=(1-\pi)\bigl|\E_Q[(\hat o_k-o)(m_0-\hat m_k)]\bigr|\le(1-\pi)\,r_o\,r_m=o_P(n^{-1/2}).
\]
Sumando sobre los bloques con pesos $n_k/n$, $P_n\psi(\cdot;\hat o,\hat m)=P_n\psi(\cdot;o,m_0)+o_P(n^{-1/2})=\pi P_n\phi+o_P(n^{-1/2})$.

\paragraph{Paso 5 (normalizaciones).} Descomponemos los factores del paso 2 en una parte i.i.d. y una de error de los parámetros auxiliares:
\begin{itemize}
  \item $\hat A=P_n[(1-Z)o(Y-m_0)]+P_n[(1-Z)\{(\hat o-o)(Y-m_0)+\hat o(m_0-\hat m)\}]$. El primer sumando es $O_P(n^{-1/2})$ por el teorema central del límite. En el segundo, la parte centrada es $o_P(n^{-1/2})$ con el argumento del paso 3, y su media condicional, por (I2), es $(1-\pi)\E_Q[\hat o_k(m_0-\hat m_k)]$, acotada por $(1-\pi)o_{\max}\,r_m$. Luego $\hat A=O_P(n^{-1/2})+O_P(r_m)$.
  \item Del mismo modo, $\hat B-\pi=(P_n-\E)[(1-Z)o]+P_n[(1-Z)(\hat o-o)]=O_P(n^{-1/2})+O_P(r_o)$, usando $\E[(1-Z)o]=\pi$ por (I1) y que la norma de $L^1(Q)$ no supera a la de $L^2(Q)$.
  \item $P_n[Z(\hat m-\theta)]=P_n[Z(m_0-\theta)]+P_n[Z(\hat m-m_0)]=O_P(n^{-1/2})+O_P(r_m)$, usando $\E[Z(m_0-\theta)]=0$ y (I3). Además, $1/\hat\pi-1/\pi=O_P(n^{-1/2})$.
\end{itemize}
Como $\hat B\to\pi>0$ en probabilidad,
\[
\hat A\Bigl(\frac1{\hat B}-\frac1\pi\Bigr)=\frac{\hat A(\pi-\hat B)}{\pi\hat B}=O_P\bigl(n^{-1}+n^{-1/2}(r_o+r_m)+r_or_m\bigr)=o_P(n^{-1/2}),
\]
y el tercer término del paso 2 es $O_P(n^{-1/2})\cdot[O_P(n^{-1/2})+O_P(r_m)]=o_P(n^{-1/2})$.

\paragraph{Paso 6 (conclusión).} $\hat\theta-\theta=P_n\phi+o_P(n^{-1/2})$. Las $\phi(O_i)$ son i.i.d., con $\E\phi=0$ y $\E\phi^2=\mathcal V<\infty$, porque $\phi$ está acotada. El teorema central del límite y el lema de Slutsky dan $\sqrt n(\hat\theta-\theta)\to N(0,\mathcal V)$.

\paragraph{Paso 7 (estimación de $\mathcal V$).} El estimador implementado usa
\[
\hat\phi(O)=\frac{(1-Z)\,\hat o(X)\,\bigl(Y-\hat m(X)-\hat A/\hat B\bigr)}{\hat B}+\frac{Z\,(\hat m(X)-\hat\mu_1)}{\hat\pi},\qquad\hat\mu_1=\frac{P_n[Z\hat m]}{\hat\pi}.
\]
Todas las funciones están acotadas, así que $|\hat\phi-\phi|\le C_0\,(|\hat o-o|+|\hat m-m_0|+|\hat\mu_1-\theta|+|\hat A/\hat B|+|\hat B-\pi|+|\hat\pi-\pi|)$ para una constante $C_0$ (aleatoria, pero acotada en ese suceso). Las medias condicionales de los cuadrados de los dos primeros términos tienden a 0 por bloques, y los escalares tienden a 0 en probabilidad, así que $P_n(\hat\phi-\phi)^2=o_P(1)$. Con $P_n\phi^2\to \mathcal V$ por la ley de los grandes números y Cauchy--Schwarz, $|P_n\hat\phi^2-P_n\phi^2|\le P_n(\hat\phi-\phi)^2+2[P_n(\hat\phi-\phi)^2\cdot P_n\phi^2]^{1/2}\to0$. Así $\hat{\mathcal V}\to \mathcal V$ en probabilidad, y por Slutsky el intervalo $\hat\theta\pm z_{1-\alpha/2}(\hat{\mathcal V}/n)^{1/2}$ tiene cobertura asintótica $1-\alpha$.

\paragraph{Varias observaciones por unidad.} Si cada unidad $i$ aporta un bloque de $N_i\le N_{\max}$ observaciones con el mismo $Z_i$, las unidades son i.i.d. y $P$ y $Q$ se definen como leyes a nivel de observación, todo lo anterior se repite con cantidades de bloque. Las identidades (I1)--(I3) valen para sumas de bloque por linealidad, con $\pi$ y $1-\pi$ sustituidos por $\bar N_1=\E[\sum_jZ_{ij}]$ y $\bar N_0=\E[\sum_j(1-Z_{ij})]$, los números esperados de observaciones objetivo y donantes por unidad (que pueden superar 1), y $o$ por los odds a nivel de observación, de modo que $o=\bar N_1w/\bar N_0$. El \emph{cross-fitting} por unidad mantiene la independencia entre el bloque evaluado y los datos de entrenamiento, y $N_i\le N_{\max}$ mantiene todo acotado. La función de influencia por unidad es $\sum_j\psi(O_{ij};o,m_0)/\bar N_1$, y $\sqrt n$ se refiere al número de unidades; el error típico agrupado es su versión muestral. Con dependencia dentro de las unidades, $\mathcal V$ es la varianza de esta función de influencia por bloques, no la forma cerrada de la sección~\ref{sec:inferencia}, que corresponde al caso i.i.d.

\section{Pesos con corte conocido en los demás escenarios}\label{app:extra}

\begin{table}[ht]
\centering\small
\caption{Pesos con corte conocido en los demás escenarios (200 réplicas, mismos datos que la tabla~\ref{tab:principal}). En cada celda, estándar $\to$ corte conocido. Recorte del Hájek en el percentil 90 calculado con $m_0$, en milésimas.}\label{tab:corte-escenarios}
\begin{tabular}{@{}lccc@{}}
\toprule
Escenario & DE del DR & Cobertura & Recorte p90 \\
\midrule
Con $C$ & 0,107 $\to$ 0,096 & 96,5 $\to$ 93,5\,\% & $+16,4\to-0,4$ \\
$C$ confunde, sin $C$ & 0,132 $\to$ 0,115 & 15,5 $\to$ 9,0\,\% & $+20,1\to-7,4$ \\
$C$ confunde, con $C$ & 0,129 $\to$ 0,113 & 17,5 $\to$ 6,5\,\% & $+32,0\to+3,9$ \\
Solapamiento medio & 0,136 $\to$ 0,147 & 94,5 $\to$ 93,5\,\% & $+19,6\to-15,6$ \\
Solapamiento bajo & 0,300 $\to$ 0,342 & 91,0 $\to$ 87,0\,\% & $+34,5\to-20,6$ \\
Confusión no observada & 0,127 $\to$ 0,109 & 0,0 $\to$ 0,0\,\% & $+23,1\to+4,6$ \\
\bottomrule
\end{tabular}
\end{table}

En los demás escenarios, con una segunda predicción registrada (tabla~\ref{tab:corte-escenarios}), los estimadores coinciden con los mismos límites que el estándar ($|z|\le1,7$), porque solo cambia cómo se estiman los mismos pesos; con confusión, la cobertura respecto de $\theta$ es incluso menor, porque el estimador es más preciso alrededor de un valor equivocado. Con la variable de canal no hay explosión de pesos: los donantes nunca la tienen registrada, y la falta de soporte aparece como $\hat e>0,99$ en el 33,3\,\% del grupo objetivo, frente a $\kappa=0,335$. Con solapamiento bajo, el recorte con estos pesos reproduce lo que la Proposición~\ref{prop:recorte} predice para los pesos verdaderos ($-2,2$, $-12,1$ y $-20,6$ frente a $-3,0$, $-11,9$ y $-20,8$). El efecto del recorte se reduce mucho en todos los escenarios y pasa a ser negativo o casi nulo, salvo en «$C$ confunde, con $C$» y en «Confusión no observada», donde sigue siendo positivo.

Fallan dos predicciones: la dispersión no baja con solapamiento medio, y con solapamiento bajo sube y la cobertura cae del 91 al 87\,\%. Nuestra interpretación, sin contrastar: al empeorar el solapamiento, la propensión de canal se acerca a 1 dentro del grupo objetivo y sus odds se estiman con pocos donantes, mientras que el clasificador estándar, con muy pocos casos positivos, los encoge hacia valores moderados. Recortar dentro del DR iguala ambos métodos: con solapamiento bajo, la raíz del error cuadrático medio en el percentil 90 es 0,160 con corte conocido y 0,163 con el estándar, con coberturas del 97,5 y el 98\,\% usando el error típico sin recortar. La ventaja en precisión se limita, por tanto, a solapamientos buenos, sobre todo con muestras pequeñas.


\begin{thebibliography}{99}\small

\bibitem[Baitairian et~al.(2025)]{baitairian2025} Baitairian, J.-B., Sebastien, B., Jreich, R., Katsahian, S. y Guilloux, A. (2025). Calibrating confounding strength in sensitivity models for weighting estimators: a comparative review and a new method. arXiv:2510.16560.

\bibitem[Bickel et~al.(2009)]{bickel2009} Bickel, S., Brückner, M. y Scheffer, T. (2009). Discriminative learning under covariate shift. \emph{Journal of Machine Learning Research}, 10, 2137--2155.

\bibitem[Blake et~al.(2020)]{blake2020} Blake, H. A., Leyrat, C., Mansfield, K. E., Seaman, S., Tomlinson, L. A., Carpenter, J. y Williamson, E. J. (2020). Propensity scores using missingness pattern information: a practical guide. \emph{Statistics in Medicine}, 39, 1641--1657.

\bibitem[Chernozhukov et~al.(2018)]{chernozhukov2018} Chernozhukov, V., Chetverikov, D., Demirer, M., Duflo, E., Hansen, C., Newey, W. y Robins, J. (2018). Double/debiased machine learning for treatment and structural parameters. \emph{The Econometrics Journal}, 21(1), C1--C68.

\bibitem[Connor Gorber et~al.(2007)]{connorgorber2007} Connor Gorber, S., Tremblay, M., Moher, D. and Gorber, B. (2007). A comparison of direct vs. self-report measures for assessing height, weight and body mass index: a systematic review. \emph{Obesity Reviews}, 8(4), 307--326.

\bibitem[Cook et~al.(2008)]{cook2008} Cook, T. D., Shadish, W. R. y Wong, V. C. (2008). Three conditions under which experiments and observational studies produce comparable causal estimates. \emph{Journal of Policy Analysis and Management}, 27(4), 724--750.

\bibitem[Crump et~al.(2009)]{crump2009} Crump, R. K., Hotz, V. J., Imbens, G. W. y Mitnik, O. A. (2009). Dealing with limited overlap in estimation of average treatment effects. \emph{Biometrika}, 96(1), 187--199.

\bibitem[D'Agostino y Rubin(2000)]{dagostino2000} D'Agostino, R. B., Jr. y Rubin, D. B. (2000). Estimating and using propensity scores with partially missing data. \emph{Journal of the American Statistical Association}, 95, 749--759.

\bibitem[Dahabreh et~al.(2019)]{dahabreh2019} Dahabreh, I. J., Robertson, S. E., Tchetgen Tchetgen, E. J., Stuart, E. A. and Hernán, M. A. (2019). Generalizing causal inferences from individuals in randomized trials to all trial-eligible individuals. \emph{Biometrics}, 75(2), 685--694.

\bibitem[D'Amour et~al.(2021)]{damour2021} D'Amour, A., Ding, P., Feller, A., Lei, L. y Sekhon, J. (2021). Overlap in observational studies with high-dimensional covariates. \emph{Journal of Econometrics}, 221, 644--654.

\bibitem[de~Luna et~al.(2011)]{deluna2011} de Luna, X., Waernbaum, I. y Richardson, T. S. (2011). Covariate selection for the nonparametric estimation of an average treatment effect. \emph{Biometrika}, 98(4), 861--875.

\bibitem[Dehejia and Wahba(1999)]{dehejia1999} Dehejia, R. H. and Wahba, S. (1999). Causal effects in nonexperimental studies: reevaluating the evaluation of training programs. \emph{Journal of the American Statistical Association}, 94(448), 1053--1062.

\bibitem[Dong et~al.(2020)]{dong2020} Dong, J., Zhang, J. L., Zeng, S. y Li, F. (2020). Subgroup balancing propensity score. \emph{Statistical Methods in Medical Research}, 29(3), 659--676.

\bibitem[Dorn y Guo(2023)]{dorn2023} Dorn, J. y Guo, K. (2023). Sharp sensitivity analysis for inverse propensity weighting via quantile balancing. \emph{Journal of the American Statistical Association}, 118(544), 2645--2657.

\bibitem[Hahn(1998)]{hahn1998} Hahn, J. (1998). On the role of the propensity score in efficient semiparametric estimation of average treatment effects. \emph{Econometrica}, 66(2), 315--332.

\bibitem[Heckman et~al.(1997)]{heckman1997} Heckman, J. J., Ichimura, H. y Todd, P. E. (1997). Matching as an econometric evaluation estimator: evidence from evaluating a job training programme. \emph{The Review of Economic Studies}, 64(4), 605--654.

\bibitem[Hirano et~al.(2003)]{hirano2003} Hirano, K., Imbens, G. W. y Ridder, G. (2003). Efficient estimation of average treatment effects using the estimated propensity score. \emph{Econometrica}, 71(4), 1161--1189.

\bibitem[Imai and Yamamoto(2010)]{imai2010} Imai, K. and Yamamoto, T. (2010). Causal inference with differential measurement error: nonparametric identification and sensitivity analysis. \emph{American Journal of Political Science}, 54(2), 543--560.

\bibitem[Khan y Tamer(2010)]{khan2010} Khan, S. y Tamer, E. (2010). Irregular identification, support conditions, and inverse weight estimation. \emph{Econometrica}, 78(6), 2021--2042.

\bibitem[LaLonde(1986)]{lalonde1986} LaLonde, R. J. (1986). Evaluating the econometric evaluations of training programs with experimental data. \emph{The American Economic Review}, 76(4), 604--620.

\bibitem[Liu et~al.(2023)]{liu2023} Liu, Y., Li, H., Zhou, Y. y Matsouaka, R. A. (2023). Average treatment effect on the treated, under lack of positivity. arXiv:2309.01334.

\bibitem[Ma y Wang(2020)]{ma2020} Ma, X. y Wang, J. (2020). Robust inference using inverse probability weighting. \emph{Journal of the American Statistical Association}, 115, 1851--1860.

\bibitem[Ma y Namkoong(2025)]{manamkoong2025} Ma, Y. y Namkoong, H. (2025). A sensitivity approach to causal inference under limited overlap. arXiv:2511.22003.

\bibitem[Ma et~al.(2023)]{ma2023} Ma, Y., Sant'Anna, P. H. C., Sasaki, Y. y Ura, T. (2023). Doubly robust estimators with weak overlap. arXiv:2304.08974.

\bibitem[McClean et~al.(2024)]{mcclean2024} McClean, A., Branson, Z. y Kennedy, E. H. (2024). Calibrated sensitivity models. arXiv:2405.08738.

\bibitem[National Center for Health Statistics(2020)]{nchs2020} National Center for Health Statistics (2020). \emph{National Health and Nutrition Examination Survey, 2017--2018 data files}. Hyattsville, MD: Centers for Disease Control and Prevention.

\bibitem[Petersen et~al.(2012)]{petersen2012} Petersen, M. L., Porter, K. E., Gruber, S., Wang, Y. y van der Laan, M. J. (2012). Diagnosing and responding to violations in the positivity assumption. \emph{Statistical Methods in Medical Research}, 21(1), 31--54.

\bibitem[Qin(1998)]{qin1998} Qin, J. (1998). Inferences for case-control and semiparametric two-sample density ratio models. \emph{Biometrika}, 85(3), 619--630.

\bibitem[Rassen et~al.(2012)]{rassen2012} Rassen, J. A., Glynn, R. J., Rothman, K., Setoguchi, S. y Schneeweiss, S. (2012). Applying propensity scores estimated in a full cohort to adjust for confounding in subgroup analyses. \emph{Pharmacoepidemiology and Drug Safety}, 21(7), 697--709.

\bibitem[Rudin(1987)]{rudin1987} Rudin, W. (1987). \emph{Real and Complex Analysis} (3.ª ed.). McGraw-Hill.

\bibitem[Smith and Todd(2005)]{smith2005} Smith, J. A. and Todd, P. E. (2005). Does matching overcome LaLonde's critique of nonexperimental estimators? \emph{Journal of Econometrics}, 125(1--2), 305--353.

\bibitem[Steiner et~al.(2011)]{steiner2011} Steiner, P. M., Cook, T. D. y Shadish, W. R. (2011). On the importance of reliable covariate measurement in selection bias adjustments using propensity scores. \emph{Journal of Educational and Behavioral Statistics}, 36(2), 213--236.

\bibitem[Sugiyama et~al.(2012)]{sugiyama2012} Sugiyama, M., Suzuki, T. y Kanamori, T. (2012). \emph{Density Ratio Estimation in Machine Learning}. Cambridge University Press.

\bibitem[Tan(2006)]{tan2006} Tan, Z. (2006). A distributional approach for causal inference using propensity scores. \emph{Journal of the American Statistical Association}, 101(476), 1619--1634.

\bibitem[VanderWeele y Shpitser(2011)]{vanderweele2011} VanderWeele, T. J. y Shpitser, I. (2011). A new criterion for confounder selection. \emph{Biometrics}, 67(4), 1406--1413.

\bibitem[Westreich et~al.(2011)]{westreich2011} Westreich, D., Cole, S. R., Funk, M. J., Brookhart, M. A. y Stürmer, T. (2011). The role of the c-statistic in variable selection for propensity score models. \emph{Pharmacoepidemiology and Drug Safety}, 20(3), 317--320.

\bibitem[Westreich et~al.(2017)]{westreich2017} Westreich, D., Edwards, J. K., Lesko, C. R., Stuart, E. and Cole, S. R. (2017). Transportability of trial results using inverse odds of sampling weights. \emph{American Journal of Epidemiology}, 186(8), 1010--1014.

\bibitem[Zhao et~al.(2019)]{zhao2019} Zhao, Q., Small, D. S. and Bhattacharya, B. B. (2019). Sensitivity analysis for inverse probability weighting estimators via the percentile bootstrap. \emph{Journal of the Royal Statistical Society: Series B}, 81(4), 735--761.

\end{thebibliography}


\end{document}
